\chapter{Basics of Ricci flow}
\label{chap:flowbasics}
\label{flowbasics}
The Ricci flow equation and its basic theory are as in
\cite{Hamilton3MPRC,ChowKnopf}.
\section{The definition of the Ricci flow}
\quad
\begin{definition}[Ricci flow equation]\label{def:ricci-flow}
\uses{def:riemannian-metric}
\lean{MorganTianLib.IsRicciFlowOn, MorganTianLib.IsInitialTime, MorganTianLib.IsInitialMetric}
\leanok

The {\em Ricci flow equation} is the
following evolution equation for a Riemannian metric:
\begin{equation}\label{Ricciflow}
\frac{\partial g}{\partial t}= -2\Ric(g). \end{equation} A solution to
this equation (or a {\em Ricci flow}) is a one-parameter
family of metrics $g(t)$, parameterized by $t$ in a non-degenerate interval
$I$, on a smooth manifold $M$ satisfying Equation~(\ref{Ricciflow}). If $I$ has
an initial point $t_0$ then $(M,g(t_0))$ is called  {\sl the initial condition
of} or {\em the initial metric for}  the
Ricci flow (or of the solution).
\end{definition}



In harmonic coordinates $(x^1,\ldots,x^n)$ about $p$,
 that is to say coordinates where $\triangle x^i=0$ for all $i$,
we have
\[ \Ric_{ij}=\Ric(\frac{\partial}{\partial x^i},\frac{\partial}{\partial x^j})=
-\frac{1}{2}\triangle g_{ij} + Q_{ij}(g^{-1},\partial g)\]
where $Q$ is quadratic in $g^{-1}$ and $\partial g$; this local form is
recorded in Lemma 3.32 of \cite{ChowKnopf}. Thus the equation is a heat equation for $g$,
\[ \frac{\partial}{\partial t}g =
\triangle g +2Q(g^{-1},\partial g). \]

\begin{definition}[Space-time]\label{def:space-time}
\label{defnRicci}
\uses{def:ricci-flow}
\lean{MorganTianLib.RicciFlowSpaceTime, MorganTianLib.RicciFlowTimeSlice, MorganTianLib.timeSlice, MorganTianLib.timeSliceEquiv, MorganTianLib.timeSliceEmbedding, MorganTianLib.timeSliceEmbedding_contMDiff, MorganTianLib.timeSliceEmbedding_isEmbedding, MorganTianLib.range_timeSliceEmbedding, MorganTianLib.timeSliceMetric, MorganTianLib.HorizontalTangentSpace, MorganTianLib.horizontalTangentInclusionAt, MorganTianLib.horizontalTangentInclusion, MorganTianLib.mfderiv_timeSliceEmbedding, MorganTianLib.horizontalTangentInclusion_injective, MorganTianLib.horizontalTangent_finrank, MorganTianLib.horizontalProjection, MorganTianLib.horizontalProjection_contMDiff, MorganTianLib.range_horizontalProjectionAt, MorganTianLib.range_horizontalProjection, MorganTianLib.horizontalMetric, MorganTianLib.horizontalMetric_timeSlice, MorganTianLib.IsRicciFlowOn.horizontalMetric_smooth, MorganTianLib.timeSliceBall, MorganTianLib.timeSliceBallInSpaceTime, MorganTianLib.backwardParabolicNeighborhood, MorganTianLib.forwardParabolicNeighborhood}
\leanok
We introduce some notation that will be used throughout. Given a
Ricci flow $(M^n,g(t))$ defined for $t$ contained in an interval
$I$, then the {\sl space-time} for this flow is  $M\times I$.  The
{\sl $t$ time-slice} of space-time is the Riemannian manifold
$M\times\{t\}$ with the Riemannian metric $g(t)$. Let ${\mathcal
HT}(M\times I)$ be the {\em horizontal tangent
bundle} of space-time, i.e., the
bundle of tangent vectors to the time-slices. It is a smooth,
rank-$n$ subbundle the tangent bundle of space-time. The evolving
metric $g(t)$ is then a smooth section of $\Sym^2{\mathcal
HT}^*(M\times I)$. We denote points of space-time as pairs $(p,t)$.
Given $(p,t)$ and any $r>0$ we denote by $B(p,t,r)$ the metric ball
of radius $r$ centered at $(p,t)$ in the $t$ time-slice. For any
$\Delta t>0$ for which $[t-\Delta t,t]\subset I$, we define the {\sl
backwards parabolic neighborhood}
$P(x,t,r,-\Delta t)$ to be the product $B(x,t,r)\times
[t-\Delta t,t]$ in space-time. Notice that the intersection of
$P(x,t,r,-\Delta t)$ with a time-slice other that the $t$ time-slice
need not be a metric ball in that time-slice. There is the
corresponding notion of a forward parabolic neighborhood
$P(x,t,r,\Delta t)$ provided that $[t,t+\Delta t]\subset I$.
\end{definition}






\section{Some exact solutions to the Ricci flow}

\subsection{Einstein manifolds}
Let $g_0$ be an Einstein metric: $\Ric(g_0) = \lambda g_0$, where $\lambda$ is a constant. Then for
any positive constant $c$, setting $g = cg_0$ we have $\Ric(g)
= \Ric(g_0) = \lambda g_0= \frac{\lambda}{c} g.$ Using this we
can construct solutions to the Ricci flow equation as follows.
Consider $g(t) = u(t)g_{0}$. If this one-parameter family of
metrics is a solution of the Ricci flow, then
\begin{align*}
\frac{\partial g}{\partial t}& =u'(t)g_{0}\\& = -2\Ric(u(t)g_{0})\\&
= -2\Ric(g_{0})\\& = -2\lambda g_{0}.
\end{align*}
So $u'(t)= -2\lambda$,  and hence $u(t) = 1-2 \lambda t$. Thus
$g(t)= (1-2 \lambda t)g_{0}$ is a solution of the Ricci flow. The
cases $\lambda
>0, \lambda=0,$ and $\lambda< 0$ correspond to \textit{shrinking}, \textit{steady}
and \textit{expanding} solutions. Notice that in the shrinking case the
solution exists for $t\in [0,\frac{1}{2\lambda})$ and goes singular at
$t=\frac{1}{2\lambda}$.

\begin{example}[Einstein metrics ($S^n, \mathbb{R}^n, \mathbb{H}^n$)]\label{ex:einstein-ricci-flow}
\uses{def:einstein-manifold}
\lean{MorganTianLib.IsEinsteinTensor, MorganTianLib.einsteinScale, MorganTianLib.einsteinScale_pos_iff_of_pos, MorganTianLib.einsteinScale_zero_lambda, MorganTianLib.einsteinScale_pos_of_neg_of_nonneg, MorganTianLib.hasDerivAt_einsteinScale, MorganTianLib.ricciTensorAt_rescaledMetric_eq, MorganTianLib.isRicciFlowEquationOn_einsteinScale, MorganTianLib.einsteinScaleExtension, MorganTianLib.einsteinMetricFamilyOn, MorganTianLib.isRicciFlowEquationOn_einsteinMetricFamilyOn, MorganTianLib.einsteinMetricFamilyOn_canonicalDist, MorganTianLib.einsteinMetricFamilyOn_sectionalCurvatureAt, MorganTianLib.isEinsteinTensor_of_isConstantCurvature, MorganTianLib.euclideanMetric_isEinsteinTensor, MorganTianLib.euclideanMetric_isRicciFlowEquationOn, MorganTianLib.hyperbolicMetric_isEinsteinTensor, MorganTianLib.hyperbolicMetric_isRicciFlowEquationOn, MorganTianLib.hyperbolicMetric_isRicciFlowEquationOn_nonneg, MorganTianLib.hyperbolicEinsteinScale_tendsto_atTop, MorganTianLib.hyperbolicMetric_canonicalDist_tendsto_atTop, MorganTianLib.hyperbolicMetric_sectionalCurvature_tendsto_zero, MorganTianLib.shrinkingRoundSphereScale, MorganTianLib.shrinkingRoundSphereScale_pos_iff, MorganTianLib.hasDerivAt_shrinkingRoundSphereScale, MorganTianLib.shrinkingRoundSphereMetricFamily, MorganTianLib.shrinkingRoundSphereMetricFamily_metricInner, MorganTianLib.shrinkingRoundSphereMetricFamily_canonicalDist, MorganTianLib.shrinkingRoundSphereMetricFamily_sectionalCurvatureAt, MorganTianLib.spherePatch_sectionalCurvature_one, MorganTianLib.spherePatch_shapeOperatorAt_identity}

The standard metric on each of $S^{n}, \Ar^{n},$ and ${\mathbb
H}^{n}$ is Einstein. Ricci flow is contracting on $S^n$, constant on
$\Ar^n$, and expanding on ${\mathbb H}^n$. The Ricci flow on $S^n$
has a finite-time singularity where the diameter of the manifold
goes to zero and the curvature goes uniformly to $+\infty$. The
Ricci flow on ${\mathbb H}^n$ exists for all $t\ge 0$ and as $t$
goes to infinity the distance between any pair of points grows
without bound and the curvature goes uniformly to zero.
\end{example}

\begin{example}[Fubini-Study metric on $\mathbb{CP}^n$]\label{ex:fubini-study-metric}
\uses{def:einstein-manifold}
\lean{MorganTianLib.HopfSphere, MorganTianLib.fact_finrank_hopfAmbient, MorganTianLib.instMulActionCircleHopfSphere, MorganTianLib.hopfSphereMetric, MorganTianLib.IsRiemannianSubmersion, MorganTianLib.FubiniStudyHopfConstruction, MorganTianLib.FubiniStudyHopfConstruction.ofExistsUnique, MorganTianLib.FubiniStudyHopfConstruction.contMDiff_projection, MorganTianLib.FubiniStudyHopfConstruction.surjective_mfderiv_projection, MorganTianLib.FubiniStudyHopfConstruction.IsHorizontal, MorganTianLib.FubiniStudyHopfConstruction.metricInner_mfderiv_projection}

$\mathbb{C}P^{n}$ equipped with the Fubini-Study metric, which is induced
from the standard metric of $S^{2n+1}$ under the Hopf fibration with the fibers
of great circles, is Einstein.
\end{example}


\begin{example}[Standard shrinking round cylinder]\label{ex:shrinking-round-cylinder}
\uses{def:ricci-flow}
\lean{MorganTianLib.shrinkingRoundSphereScale, MorganTianLib.hasDerivAt_shrinkingRoundSphereScale, MorganTianLib.shrinkingRoundSphereScale_tendsto_zero_at_extinction, MorganTianLib.shrinkingRoundSphereScale_tendsto_zero_atTop, MorganTianLib.shrinkingRoundSphereMetricFamily, MorganTianLib.shrinkingRoundSphereMetricFamily_metricInner, MorganTianLib.shrinkingRoundSphereMetricFamily_canonicalDist, MorganTianLib.shrinkingRoundSphereMetricFamily_canonicalDist_tendsto_zero_atTop, MorganTianLib.shrinkingRoundSphereMetricFamily_sectionalCurvatureAt}

Let $h_0$ be the round metric on $S^2$ with constant Gausssian
curvature $1/2$. Set $h(t)=(1-t)h_0$. Then the flow $$(S^2,h(t)),\
-\infty<t<1,$$ is a Ricci flow. We also have the product of this flow
with the trivial flow on the line: $(S^2\times \Ar,h(t)\times
ds^2),\ -\infty<t<1$. This is called the {\em standard shrinking
round cylinder}.
\end{example}

The evolving neck definition is the backward-in-time version of an
$\epsilon$-neck used for flow limits.


\begin{definition}[Evolving $\epsilon$-neck]\label{def:evolving-epsilon-neck}
\uses{def:epsilon-neck}
\label{strongepsneck}
\lean{MorganTianLib.evolvingEpsilonNeckTimeSet, MorganTianLib.evolvingEpsilonNeckActualTime, MorganTianLib.IsStandardShrinkingCylinderFamily, MorganTianLib.EvolvingEpsilonNeck, MorganTianLib.StrongEpsilonNeck, MorganTianLib.evolvingEpsilonNeckTimeSet_nonempty_iff, MorganTianLib.evolvingEpsilonNeckZeroTime_value, MorganTianLib.evolvingEpsilonNeckActualTime_zeroTime, MorganTianLib.evolvingEpsilonNeckActualTime_rescaled_sub, MorganTianLib.evolvingEpsilonNeckActualTime_mem_Ioc, MorganTianLib.evolvingEpsilonNeckActualTime_image_timeSet, MorganTianLib.evolvingEpsilonNeckActualTime_strictMono, MorganTianLib.EvolvingEpsilonNeck.actualTime_mem_domain, MorganTianLib.EvolvingEpsilonNeck.centralTime_eq_t0, MorganTianLib.standardShrinkingCylinder_scale_pos_of_mem, MorganTianLib.standardShrinkingCylinder_scale_ge_one_of_mem, MorganTianLib.standardShrinkingCylinder_scalarCurvature_pos_of_mem, MorganTianLib.standardShrinkingCylinderInner_zero, MorganTianLib.standardShrinkingCylinderFamily_metricInner_zero}
\leanok

Let $(M,g(t))$ be a Ricci flow. An {\em evolving $\epsilon$-neck
centered at $(x,t_0)$ and defined for rescaled time
$t_1$} is an $\epsilon$-neck
$$\varphi\colon S^2\times
(-\epsilon^{-1},\epsilon^{-1})\xrightarrow{\cong}
N\subset (M,g(t))$$ centered at $(x,t_0)$ with the property that
pull-back via $\varphi$ of the family of metrics $R(x,t_0)g(t')|_N,\
-t_1< t'\le 0$, where $t_1=R(x,t_0)^{-1}(t-t_0)$, is within
$\epsilon$ in the $C^{[1/\epsilon]}$-topology of the product of the
standard metric on the interval with evolving round metric on $S^2$
with scalar curvature $1/(1-t')$ at time $t'$. A {\em strong
$\epsilon$-neck centered at
$(x,t_0)$} in a Ricci flow is an
evolving $\epsilon$-neck centered at $(x,t_0)$ and defined for
rescaled time $1$, see the figure in [Morgan--Tian].
\end{definition}


\subsection{Solitons}

A {\sl Ricci soliton}
is a Ricci flow $(M,g(t)),\ 0\le t<T\le \infty$, with the property that for
each $t\in [0,T)$ there is a diffeomorphism $\varphi_t\colon M\to M$ and a
constant $\sigma(t)$ such that $\sigma(t)\varphi_t^*g(0)=g(t)$. That is to say,
in a Ricci soliton all the Riemannian manifolds $(M,g(t))$ are isometric up to
a scale factor that is allowed to vary with $t$. The soliton is said to be {\em
shrinking} if $\sigma'(t)<0$ for all $t$. One way to generate Ricci solitons is
the following: Suppose that we have a vector field $X$ on $M$ and a constant
$\lambda$ and a metric $g(0)$ such that
\begin{equation}\label{soliton}
-\Ric(g(0))=\frac{1}{2}{\mathcal L}_Xg(0)-\lambda g(0).
\end{equation} We set $T=\infty$ if $\lambda\le 0$ and equal to $(2\lambda)^{-1}$
if $\lambda>0$. Then, for all $t\in [0,T)$ we define a function
$$\sigma(t)=1-2\lambda t,$$
and a vector field
$$Y_t(x)=\frac{X(x)}{\sigma(t)}.$$
 Then we define $\varphi_t$ as the one-parameter
family of diffeomorphisms generated by the time-dependent vector
fields $Y_t$.
\begin{claim}[Soliton generation]\label{claim:soliton-generation}
\uses{def:ricci-flow}
\label{solitonclaim}
\lean{MorganTianLib.RicciSoliton, MorganTianLib.RicciSoliton.IsShrinking, MorganTianLib.metricLieDerivativeAt, MorganTianLib.IsSolitonGenerator, MorganTianLib.solitonScale, MorganTianLib.solitonTimeSet, MorganTianLib.hasDerivAt_solitonScale, MorganTianLib.solitonScale_pos_of_mem, MorganTianLib.solitonScale_isShrinking, MorganTianLib.solitonFlowTime, MorganTianLib.solitonFlowTime_zero, MorganTianLib.hasDerivAt_solitonFlowTime, MorganTianLib.hasDerivAt_solitonFlowTime_of_mem, MorganTianLib.mem_solitonTimeSet_iff_scale_pos_of_nonneg, MorganTianLib.solitonScale_ge_one_of_nonpos, MorganTianLib.solitonScale_at_shrinking_endpoint, MorganTianLib.solitonScale_eq_endpoint_gap, MorganTianLib.solitonScale_strictAnti, MorganTianLib.solitonScale_strictMono, MorganTianLib.solitonScale_unbounded_of_neg, MorganTianLib.solitonTimeSet_eq_Ico_of_pos, MorganTianLib.solitonTimeSet_eq_Ici_of_nonpos, MorganTianLib.exp_neg_two_mul_solitonFlowTime_eq_scale, MorganTianLib.solitonFlowTime_nonneg_of_mem, MorganTianLib.hasDerivAt_solitonScale_inv, MorganTianLib.solitonScale_mul_solitonFlowTime_deriv_of_mem}

 The flow
$(M,g(t)),\ 0\le t<T$, where $g(t)=\sigma(t)\varphi_t^*g(0)$, is a soliton. It
is a shrinking soliton if $\lambda>0$.
\end{claim}

\begin{proof}
\sketch Verify the Ricci flow equation:
that, the result follows immediately. We have
\begin{eqnarray*} \frac{\partial g(t)}{\partial t} & = &
\sigma'(t)\varphi_t^*g(0)+\sigma(t)\varphi_t^*{\mathcal L}_{Y(t)}g(0)\\
& = & \varphi_t^*(-2\lambda+{\mathcal L}_X)g(0)\\
& = & \varphi_t^*(-2\Ric(g(0))) =-2\Ric(\varphi_t^*(g(0))).
\end{eqnarray*}
Since $\Ric(\alpha g)=\Ric(g)$ for any $\alpha>0$, it follows that
$$\frac{\partial g(t)}{\partial t} = -2\Ric(g(t)).$$
\end{proof}



\begin{definition}[Gradient shrinking soliton]\label{def:gradient-shrinking-soliton}
\lean{MorganTianLib.GradientShrinkingRicciSoliton}
\leanok
A shrinking soliton $(M,g(t)),\ 0\le t<T$, is said to be a {\em
gradient shrinking soliton} if
the vector field $X$ in Equation~(\ref{soliton}) is the gradient of
a smooth function $f$ on $M$.
\end{definition}

\begin{proposition}[Gradient shrinking soliton from Hessian]\label{prop:gradient-shrinking-soliton-generation}
\uses{def:gradient-shrinking-soliton, lem:hessian-symmetric}
\label{GSSgeneration}
\lean{MorganTianLib.metricLieDerivativeAt_gradientField, MorganTianLib.IsGradientShrinkerPotential, MorganTianLib.IsGradientShrinkerPotential.isSolitonGenerator}

 Suppose we have a complete Riemannian manifold $(M,g(0))$, a
smooth function $f\colon M\to \Ar$, and a constant $\lambda>0$ such
that
\begin{equation}\label{GSS}
-\Ric(g(0))=\Hess(f)-\lambda g(0). \end{equation}
 Then there is $T>0$ and a gradient shrinking soliton
$(M,g(t))$ defined for  $0\le t<T.$ \end{proposition}

\begin{proof}
\sketch Since
$${\mathcal L}_{\nabla f}g(0)=2\Hess(f),$$
 Equation~(\ref{GSS}) is the soliton equation,
Equation~(\ref{soliton}), with the vector field $X$ being the
gradient vector field $\nabla f$. It is a shrinking soliton by
assumption since $\lambda>0$.
\end{proof}

\begin{definition}[Generating gradient shrinking soliton]\label{def:gradient-shrinking-soliton-generator}
\uses{def:gradient-shrinking-soliton}
\lean{MorganTianLib.GeneratesGradientShrinkingSoliton}
\leanok

In this case we say that $(M,g(0))$ and $f\colon M\to \Ar$ {\em
generate} a gradient shrinking soliton.
\end{definition}


\section{Local existence and uniqueness}\label{3.3}

\begin{theorem}[Local existence and uniqueness (Hamilton)]\label{thm:local-existence-uniqueness}
\uses{def:ricci-flow}
\lean{MorganTianLib.RicciCovector, MorganTianLib.ricciCovectorNormSq, MorganTianLib.ricciMatrixNormSq, MorganTianLib.ricciMatrixPairing, MorganTianLib.ricciCovectorNormSq_nonneg, MorganTianLib.ricciCovectorNormSq_pos, MorganTianLib.ricciMatrixNormSq_nonneg, MorganTianLib.ricciMatrixNormSq_pos, MorganTianLib.ricciLinearisationSymbol, MorganTianLib.ricciGaugeSymbol, MorganTianLib.deTurckGaugeCancellationSymbol, MorganTianLib.deTurckLinearisationSymbol, MorganTianLib.ricciLinearisationSymbol_rankOne, MorganTianLib.ricciLinearisationSymbol_gaugeKernel, MorganTianLib.deTurckLinearisationSymbol_eq_scalar, MorganTianLib.deTurckLinearisationSymbol_pairing_self, MorganTianLib.deTurckLinearisationSymbol_strictlyParabolic, MorganTianLib.deTurckLinearisationSymbol_coercive, MorganTianLib.metricRiesz_eq_chartInvGram_sum, MorganTianLib.ricciDeTurckChartQuadratic_cotangent, MorganTianLib.ricciDeTurckChartQuadratic_cotangent_pos, MorganTianLib.ricciDeTurckChartSymbol_pairing_self, MorganTianLib.ricciDeTurckChartSymbol_coercive, MorganTianLib.RicciDeTurckChartStrictParabolic, MorganTianLib.canonicalRicciDeTurckChartStrictParabolic, MorganTianLib.ricciDeTurckVariation, MorganTianLib.IsRicciDeTurckEquationOn, MorganTianLib.IsGaugePullbackOn, MorganTianLib.IsRicciPullbackCompatible, MorganTianLib.IsHamiltonGaugeTransportOn, MorganTianLib.gaugeTransportVariationValue_eq_neg_two_pullbackRicci, MorganTianLib.isRicciFlowEquationOn_of_hamiltonGaugeTransport, MorganTianLib.isRicciFlowEquationOn_of_hamiltonGaugeTransport_predicate, MorganTianLib.MetricFamilyAgreementOn, MorganTianLib.metricFamilyAgreementOn_of_sameGauge, MorganTianLib.metricFamilyAgreementOn_of_sameGauge_predicate, MorganTianLib.metricFamily_eq_of_sameGauge, MorganTianLib.metricFamily_eq_of_sameGauge_predicate, MorganTianLib.gaugePullbackValue_refl, MorganTianLib.initialMetric_eq_of_hamiltonGaugePullback, MorganTianLib.initialMetric_eq_of_hamiltonGaugePullback_predicate, MorganTianLib.HamiltonGaugeCertificate.smooth, MorganTianLib.HamiltonGaugeCertificate.deTurckSmooth, MorganTianLib.HamiltonGaugeCertificate.deTurckEquation, MorganTianLib.HamiltonGaugeCertificate.deTurckInitial, MorganTianLib.HamiltonGaugeCertificate.pullback, MorganTianLib.HamiltonGaugeCertificate.transport, MorganTianLib.HamiltonGaugeCertificate.ricciNaturality, MorganTianLib.HamiltonGaugeCertificate.initial, MorganTianLib.HamiltonGaugeCertificate.isRicciFlowOn, MorganTianLib.exists_localRicciFlow_of_hamiltonGaugeCertificate, MorganTianLib.RicciDeTurckLocalSolution, MorganTianLib.RicciDeTurckLocalSolution.smooth, MorganTianLib.RicciDeTurckLocalSolution.deTurckEquation, MorganTianLib.RicciDeTurckLocalSolution.initial, MorganTianLib.HamiltonGaugeTransport, MorganTianLib.HamiltonGaugeTransport.certificate, MorganTianLib.exists_localRicciFlow_of_splitHamiltonGauge, MorganTianLib.HamiltonGaugeTransport.initial, MorganTianLib.HamiltonGaugeTransport.isRicciFlowOn, MorganTianLib.HamiltonGaugeTransport.eq_on_overlap, MorganTianLib.HamiltonGaugeTransport.agreement_on_overlap, MorganTianLib.RicciDeTurckStrictParabolic, MorganTianLib.canonicalRicciDeTurckStrictParabolic, MorganTianLib.RicciDeTurckClassicalOutput, MorganTianLib.RicciDeTurckClassicalOutput.toLocalSolution, MorganTianLib.RicciDeTurckPDESolver, MorganTianLib.canonicalRicciDeTurckPDESolverData, MorganTianLib.RicciDeTurckPDESolver.localSolution, MorganTianLib.exists_ricciDeTurckLocalSolution_of_solver, MorganTianLib.ricciDeTurckLocalSolution_of_solver_has_initial_and_equation, MorganTianLib.RicciDeTurckClassicalOutput.metric_eq_on_overlap, MorganTianLib.RicciDeTurckClassicalOutput.metric_eq_on_overlap_of_sharedVariation, MorganTianLib.ricciBilin_naturality_of_basis_map, MorganTianLib.gaugePullbackRicciValue_refl, MorganTianLib.isGaugePullbackOn_refl, MorganTianLib.isRicciPullbackCompatible_refl, MorganTianLib.IsRicciFlowOn.restrict, MorganTianLib.metricFamilyAgreementOn_of_sharedMetricVariation, MorganTianLib.metricFamilyAgreementOn_of_equal_ricci_on_shared_initial, MorganTianLib.metricFamily_eq_of_equal_ricci_on_shared_initial}
\lean{MorganTianLib.ParabolicPDE.CompleteInvariantContraction, MorganTianLib.ParabolicPDE.CompleteInvariantContraction.exists_fixedPoint, MorganTianLib.ParabolicPDE.CompleteInvariantContraction.fixedPoint_unique, MorganTianLib.ParabolicPDE.UniformContractionFamily, MorganTianLib.ParabolicPDE.UniformContractionFamily.fixedPointAt, MorganTianLib.ParabolicPDE.UniformContractionFamily.dist_fixedPointAt_le_of_lipschitz, MorganTianLib.ParabolicPDE.UniformContractionFamily.continuous_fixedPointAt_of_lipschitz, MorganTianLib.RicciDeTurckPicardModel, MorganTianLib.RicciDeTurckPicardModel.fixedPoint, MorganTianLib.RicciDeTurckPicardModel.classicalOutput, MorganTianLib.RicciDeTurckPicardModel.localSolution, MorganTianLib.RicciDeTurckPicardModel.localSolution_initial, MorganTianLib.RicciDeTurckPicardModel.localSolution_equation, MorganTianLib.RicciDeTurckPicardModel.exists_ricciDeTurckLocalSolution_of_picard, MorganTianLib.RicciDeTurckPicardFamily, MorganTianLib.RicciDeTurckPicardFamily.fixedPoint_parameter_stability, MorganTianLib.RicciDeTurckPicardFamily.fixedPoint_parameter_continuous, MorganTianLib.gaugePullbackMetric, MorganTianLib.gaugePullbackMetric_metricInner, MorganTianLib.gaugePullbackMetric_dcPreservesMetric, MorganTianLib.gaugePushforwardVectorField_apply, MorganTianLib.gaugePushforwardVectorField_bracket, MorganTianLib.gaugePushforwardVectorField_dir, MorganTianLib.gaugePushforward_leviCivita_cov, MorganTianLib.gaugePushforward_leviCivita_curvature, MorganTianLib.curvatureFormAt_gauge_naturality, MorganTianLib.gaugePullbackMetric_ricciTensorAt, MorganTianLib.gaugePullbackMetricFamily, MorganTianLib.gaugePullbackMetricFamily_isGaugePullbackOn, MorganTianLib.gaugePullbackMetricFamily_isRicciPullbackCompatible, MorganTianLib.gaugePullbackMetricFamily_zero, MorganTianLib.HamiltonGaugeTransport.of_concretePullback, MorganTianLib.HamiltonGaugeTransport.of_concretePullback_natural}
\lean{MorganTianLib.RicciDeTurckPDESolver.of_picard_models, MorganTianLib.exists_ricciDeTurckLocalSolution_of_picard_models, MorganTianLib.SmoothTimeDependentVectorField, MorganTianLib.TimeDependentFlowBox, MorganTianLib.exists_timeDependentFlowBox, MorganTianLib.TimeDependentFlowBox.time_coord_of_mem_slice, MorganTianLib.TimeDependentFlowBox.spatialFlow, MorganTianLib.TimeDependentFlowBox.spatialFlow_zero, MorganTianLib.TimeDependentFlowBox.spatialFlow_hasMFDerivAt, MorganTianLib.TimeDependentFlowBox.eqOn_of_common, MorganTianLib.TimeDependentFlowBox.spatialFlow_eq_of_common}
\label{compuniq}

{\bf (Hamilton, cf.
\cite{Hamilton3MPRC}.)} Let $(M,g_0)$ be a compact
Riemannian manifold of dimension $n$.
\begin{enumerate}
\item There is a $T>0$ depending on $(M,g_0)$ and  a Ricci flow
$(M,g(t)),\  0\le t<T$, with $g(0)=g_0$.
\item Suppose that we have Ricci flows with initial conditions $(M,g_0)$ at time
$0$ defined respectively on time intervals $I$ and $I'$. Then these
flows agree on $I\cap I'$. \end{enumerate}
\end{theorem}

\begin{proof}
\sketch The DeTurck argument, replacing the Nash--Moser approach
\cite{HamiltonIFTNM} by gauge breaking \cite{DeTurck}, begins with the first
variation at a Riemannian metric $g$ of minus twice the Ricci tensor in the direction $h$;
the construction is the one in \cite{DeTurck} and \cite{ChowKnopf}:
\begin{align*}
\delta_g(-2\Ric)(h)=\triangle h-\Sym(\nabla V)+S
\end{align*}
where:
\begin{enumerate}
\item $V$ is the one-form given by
$$V_k=\frac{1}{2}g^{pq}(\nabla_ph_{qk}+\nabla_qh_{pk}-\nabla_kh_{pq}),$$
\item $\Sym(\nabla V)$ is the symmetric two-tensor obtained by
symmetrizing the covariant derivative of
$V$, and
\item  $S$ is a symmetric two-tensor constructed from the inverse of the metric,
the Riemann curvature tensor and $h$, but involves no derivatives of $h$.
\end{enumerate}

Now let $g_0$ be the initial metric. For any metric $g$ we define a one-form
$\hat W$ by taking the trace, with respect to $g$, of the matrix-valued
one-form that is the difference of the connections of $g$ and $g_0$. Now we
form a second-order operator of $g$ by setting
$$P(g)={\mathcal L}_Wg,$$
the Lie derivative of $g$ with respect to the vector field $W$ dual to $\hat
W$. Thus, in local coordinates we have
$P(g)_{ij}=\nabla_i\hat W_j+\nabla_j\hat W_i$.
The linearization at $g$ of the second-order operator $P$ in the direction
$h$ is symmetric and is given by
$$\delta_gP(h)=\Sym(\nabla V)+T$$
where $T$ is a first-order operator in $h$.
Thus, defining $Q=-2\Ric+P$ we have
$$\delta_g(Q)(h)=\triangle h+U$$
where $U$ is a first-order operator in $h$.
Now we introduce the Ricci-DeTurck  flow

\begin{align}\label{brokeneqn}
\frac{\partial g}{\partial t}&=-2\Ric(g)+P.
\end{align}
The computations above show that the Ricci-DeTurck flow is  strictly parabolic.
 Thus, Equation~(\ref{brokeneqn})  has a
short-time solution $\bar g(t)$ with $\bar g(0)=g_0$ by the standard PDE theory.
Given this solution $\bar g(t)$ we
define the time-dependent vector field $W(t)=W(\bar g(t),g_0)$ as above.
 Let $\phi_t$ be a one-parameter family of diffeomorphisms, with $\phi_0=\mathrm{Id}$,
generated by this time-dependent vector field, i.e.,
$$\frac{\partial\phi_t}{\partial t}=W(t).$$
Then, direct computation shows that $g(t)=\phi_t^*\bar g(t)$ solves the Ricci flow equation.
\end{proof}



\begin{proposition}[Smooth patching of Ricci flows]\label{prop:smooth-patching}
\uses{def:ricci-flow}
\lean{MorganTianLib.patchedMetricFamily, MorganTianLib.patchedMetricFamily_hasDerivWithinAt_Iic_at_join, MorganTianLib.patchedMetricFamily_hasDerivAt_at_join, MorganTianLib.isRicciFlowEquationOn_patchedMetricFamily, MorganTianLib.isRicciFlowOn_patchedMetricFamily_of_smooth, MorganTianLib.leftEndpointLimit_unique, MorganTianLib.leftEndpoint_tendsto_of_tendstoUniformlyOn, MorganTianLib.leftEndpoint_tendsto_of_tendstoLocallyUniformlyOn, MorganTianLib.leftEndpoint_tendsto_of_continuousAt, MorganTianLib.EndpointCoefficientLimits, MorganTianLib.endpointCoefficientLimits_of_continuousAt, MorganTianLib.EndpointCoefficientLimits.patched_metric, MorganTianLib.EndpointCoefficientLimits.patched_ricci, MorganTianLib.endpointMetricCoefficient_unique_of_limits, MorganTianLib.endpointRicciCoefficient_unique_of_limits, MorganTianLib.RicciQuadraticControlOn, MorganTianLib.hasAbsoluteRicciBoundOn_of_ricciQuadraticControlOn, MorganTianLib.metricInner_exp_comparison_of_ricciQuadraticControlOn, MorganTianLib.metricInner_uniform_exp_comparison_of_ricciQuadraticControlOn, MorganTianLib.exists_metricCoefficient_limit_of_ricciQuadraticControlOn, MorganTianLib.exists_pos_metricCoefficient_limit_of_ricciQuadraticControlOn, MorganTianLib.ricciQuadraticControlOn_of_hasCurvatureOperatorNormLeOnTime, MorganTianLib.exists_metricCoefficient_limit_of_curvatureOperatorNormLeOnTime, MorganTianLib.exists_pos_metricCoefficient_limit_of_curvatureOperatorNormLeOnTime, MorganTianLib.endpointCoefficientLimits_of_supplied_endpoint_limits, MorganTianLib.EndpointFiberBilinearForm, MorganTianLib.exists_endpointFiberBilinearForm_of_ricciQuadraticControlOn, MorganTianLib.exists_endpointFiberBilinearForm_of_curvatureOperatorNormLeOnTime, MorganTianLib.RicciFlowExtensionOn, MorganTianLib.SmoothEndpointExtension, MorganTianLib.smoothMetricFamilyOn_patchedMetricFamily_of_extension, MorganTianLib.SmoothEndpointRestartData, MorganTianLib.SmoothEndpointRestartData.toAdapter, MorganTianLib.SmoothEndpointRestart, MorganTianLib.SmoothEndpointRestart.isRicciFlowOn, MorganTianLib.SmoothEndpointRestart.agrees_left, MorganTianLib.SmoothEndpointRestart.agrees_right, MorganTianLib.SmoothEndpointRestart.extends_left, MorganTianLib.SmoothEndpointRestart.extends_right, MorganTianLib.SmoothEndpointRestart.hasDerivAt_join, MorganTianLib.isRicciFlowOn_patchedMetricFamily_of_closed_left, MorganTianLib.LocalSmoothEndpointGerm, MorganTianLib.smoothMetricFamilyOn_patchedMetricFamily_of_local_germ, MorganTianLib.SmoothEndpointRestart.of_local_germ, MorganTianLib.isRicciFlowOn_patchedMetricFamily_of_local_germ, MorganTianLib.tendsto_iteratedDeriv_of_contDiffAt}
\lean{MorganTianLib.RicciFlowExtensionOn.trans, MorganTianLib.RicciFlowExtensionOn.restrict_left, MorganTianLib.SmoothEndpointRestart.extends_left_through, MorganTianLib.patchedMetricFamily_assoc, MorganTianLib.Ico_subset_Ico_of_endpoint_patch}
\lean{MorganTianLib.isMetricVariationOn_comp_sub_const, MorganTianLib.isRicciFlowEquationOn_comp_sub_const, MorganTianLib.isRicciFlowEquationOn_of_timeShift_eq_Ico, MorganTianLib.isRicciFlowOn_of_timeShift_eq_Ico_of_smooth, MorganTianLib.EndpointFiberBilinearForm.toContinuousLinearMap, MorganTianLib.smoothEndpointMetric, MorganTianLib.smoothEndpointMetric_tendsto_of_endpointField, MorganTianLib.exists_smoothEndpointMetric_of_globalEndpointField}
\label{patch}

 Suppose that $(U,g(t)),\ a\le t<b$, is a Ricci flow and suppose that
 there is a Riemannian metric $g(b)$ on $U$ such that
 as $t\rightarrow b$ the metrics $g(t)$ converge in the $C^\infty$-topology, uniformly
 on compact subsets, to $g(b)$. Suppose also that $(U,g(t)),\ b\le t<c$, is a
 Ricci flow. Then the one-parameter family of metrics $g(t), \ a\le
 t<c$, is a $C^\infty$-family and is a solution to the Ricci flow
 equation on the entire interval $[a,c)$.
 \end{proposition}

%See\note{Give reference.} ??.

\section{Evolution of curvatures}\label{evolcurv}


Let us fix a set
$(x^1,\ldots,x^n)$ of local coordinates. The Ricci flow equation,
written in local coordinates
\[ \frac{\partial g_{ij}}{\partial t} =
-2\Ric_{ij}
\]  implies a heat equation for the Riemann
curvature tensor $R_{ijkl}$ which we now derive. Various second-order
derivatives of the curvature tensor are likely to differ by
terms quadratic in the curvature tensors. To this end we introduce
the tensor $$B_{ijkl} = g^{pr}g^{qs}R_{ipjq}R_{krls}.$$ Note that we
have the obvious symmetries
$$B_{ijkl} = B_{jilk} = B_{klij},$$ but the other symmetries of
the curvature tensor $R_{ijkl}$ may fail to hold for $B_{ijkl}$.

\smallskip

\begin{theorem}[Evolution of curvatures]\label{thm:curvature-evolution}
\uses{def:ricci-flow}
\lean{MorganTianLib.curvatureB, MorganTianLib.curvatureB_swap_pairs, MorganTianLib.curvatureB_swap_within, MorganTianLib.exists_chartFrame_hasDerivAt_chartRiemannCoefOnE_curvatureEvolution_explicit_of_isRicciFlowOn, MorganTianLib.exists_chartFrame_hasDerivAt_chartRicciCoefOnE_evolution_of_isRicciFlowOn, MorganTianLib.hasDerivAt_scalarCurvatureAt_leviCivita_of_isRicciFlowOn_eq_laplacian_add_reaction, MorganTianLib.hasDerivWithinAt_chartVolumeDensity_of_isRicciFlowOn}
\leanok
\label{thmRmevol}

The curvature tensor $R_{ijkl}$, the Ricci curvature $\Ric_{ij}$, the scalar
curvature $R$, and the volume form $d\vol(x,t)$ satisfy the following evolution
equations under Ricci flow:
\begin{eqnarray}
\frac{\partial R_{ijkl}}{\partial t} &= & \triangle R_{ijkl} +2(B_{ijkl} -
B_{ijlk}-B_{iljk} + B_{ikjl})\nonumber \\&  &  - g^{pq}(R_{pjkl}\Ric_{qi} +
R_{ipkl}\Ric_{qj} + R_{ijpl}\Ric_{qk} +
R_{ijkp}\Ric_{ql}) \label{Rmevol} \\
\label{Ricevol} \frac{\partial}{\partial t}\Ric_{jk} & = & \triangle
\Ric_{jk}+2g^{pq}g^{rs}R_{jpkr}\Ric_{qs}-2g^{pq}\Ric_{jp}\Ric_{qk} \\
\label{Revol}\frac{\partial}{\partial t}R & = & \Delta R+2|\Ric|^2 \\
\label{dvolevol} \frac{\partial}{\partial t}d\vol(x,t) & = &
-R(x,t)d\vol(x,t).
\end{eqnarray}
\end{theorem}

These equations are contained in Lemma 6.15  on page 179, Lemma 6.9
 on page 176,  Lemma 6.7 on page 176, and Equation (6.5) on
page 175 of \cite{ChowKnopf}, respectively.


Let us derive some consequences of these evolution equations.
The first result is obvious from the Ricci flow equation and will be used implicitly throughout
the paper.

\begin{lemma}[Distance non-increasing under positive Ricci flow]\label{lem:distance-nonincreasing}
\uses{def:ricci-flow}
\lean{MorganTianLib.MetricInnerLE, MorganTianLib.metricInnerLE_of_isRicciFlowOn_nonnegativeRicci, MorganTianLib.metricPathIntegral_metricInnerLE, MorganTianLib.metricPathIntegral_metricInnerLE_mul, MorganTianLib.metricIntrinsicEDist_metricInnerLE, MorganTianLib.metricIntrinsicEDist_metricInnerLE_mul, MorganTianLib.distance_nonincreasing_of_isRicciFlowOn_nonnegativeRicci, MorganTianLib.metricIntrinsicEDist_eq_riemannianEDist}
\leanok
\label{posdistdec}

Suppose that $(M,g(t)),\ a<t<b$ is a Ricci flow of non-negative Ricci curvature
with $M$ a connected manifold.
Then for any points $x,y\in M$ the function $d_{g(t)}(x,y)$ is a non-increasing function of $t$.
\end{lemma}


\begin{proof}
\sketch The Ricci flow equation and non-negative Ricci curvature imply
$\partial g/\partial t\le 0$. Hence, the length of any tangent vector in $M$, and consequently the length
of any path in $M$, is a non-increasing function of $t$. Since the distance between points is the infimum
over all rectifiable paths from $x$ to $y$ of the length of the path, this function is also a non-increasing function
of $t$.
\end{proof}



\begin{lemma}[Metric and volume distortion bounds]\label{lem:metric-volume-distortion}
\uses{def:ricci-flow}
\lean{MorganTianLib.HasCurvatureOperatorNormLeOnTime, MorganTianLib.HasAbsoluteRicciBoundOn, MorganTianLib.metricInner_exp_comparison_of_isRicciFlowOn, MorganTianLib.metricInner_uniform_exp_comparison_of_isRicciFlowOn, MorganTianLib.metricInner_distortion_of_curvatureOperatorNormLeOnTime, MorganTianLib.exists_metricDistortionConstant_of_curvatureOperatorNormLeOnTime, MorganTianLib.metricIntrinsicEDist_bilipschitz_of_metricInner_bilipschitz, MorganTianLib.exists_metricDistanceDistortionConstant_of_curvatureOperatorNormLeOnTime, MorganTianLib.IsVolumeDensityEvolution, MorganTianLib.HasAbsoluteScalarBoundOn, MorganTianLib.abs_scalarCurvatureAt_le_finrank_sq_mul_of_hasCurvatureOperatorNormLeAt, MorganTianLib.HasAbsoluteScalarCurvatureBoundOnTime, MorganTianLib.hasAbsoluteScalarCurvatureBoundOnTime_of_hasCurvatureOperatorNormLeOnTime, MorganTianLib.volumeDensity_exp_comparison, MorganTianLib.chartVolumeDensity_exp_comparison_of_curvatureOperatorNormLeOnTime, MorganTianLib.realDensityMeasure, MorganTianLib.realDensityMeasure_apply, MorganTianLib.realDensityMeasure_mul_le_of_aemeasurable, MorganTianLib.realDensityMeasure_le_mul_of_aemeasurable, MorganTianLib.realDensityMeasure_exp_comparison_of_aemeasurable, MorganTianLib.realDensityMeasure_exp_comparison, MorganTianLib.realDensityMeasure_exp_comparison_of_isRealDensityEvolution, MorganTianLib.realDensityMeasure_exp_comparison_of_isOpen, MorganTianLib.chartVolumeDensity_aemeasurable_on_chartPreimage, MorganTianLib.riemannianMeasure_chartDensity_comparison, MorganTianLib.riemannianMeasure_chartDensity_comparison_of_aemeasurable, MorganTianLib.riemannianMeasure_chartDensity_exp_comparison_of_curvatureOperatorNormLeOnTime, MorganTianLib.riemannianMeasure_chartDensity_exp_comparison_of_curvatureOperatorNormLeOnTime_of_smooth}

Suppose that $(M,g(t)),\ 0\le t\le T$, is a Ricci flow and
$|Rm(x,t)|\le K$ for all $x\in M$ and all $t\in [0,T]$. Then there
are constants $A,A'$  depending on $K,T$ and the dimension such
that:
\begin{enumerate}
\item For any non-zero tangent vector $v\in T_xM$ and any $t\le T$ we have
$$A^{-1}\langle v,v\rangle_{g(0)}\le \langle v,v\rangle_{g(t)}\le A\langle
v,v\rangle_{g(0)}.$$
\item For any open subset $U\subset M$ and any $t\le T$ we have
$$(A')^{-1}\Vol_0(U)\le \Vol_t(U)\le A'\Vol_0(U).$$
\end{enumerate}
\end{lemma}

\begin{proof}
\sketch The metric and volume evolution equations give the bounds after
integrating the curvature estimate.
\uses{thm:curvature-evolution}
The Ricci flow equation yields
$$\frac{d}{dt}\left(\langle v,v\rangle_{g(t)}\right)=-2\Ric(v,v).$$
The bound on the Riemann curvature gives a bound on $\Ric$.
Integrating yields the result. The second statement is proved
analogously using Equation~(\ref{dvolevol}).
\end{proof}



\section{Curvature evolution in an evolving orthonormal
frame}\label{evolframe}


It is often best to study the evolution of the representative of the
tensor in an orthonormal frame $F$. Let $(M,g(t)),\ 0\le t<T$, be a
Ricci flow, and suppose that ${\mathcal F}$ is a frame on an open
subset $U\subset M$ consisting of vector fields $\{F_{1},
F_{2},\cdots,F_{n}\}$ on $U$ that are $g(0)$-orthonormal at every
point. Since the metric evolves by the Ricci flow, to keep the frame
orthonormal we must evolve it by an equation involving Ricci
curvature. We evolve this local frame  according to the formula
\begin{equation}\label{frameevol} \frac{\partial F_{a}}{\partial t}
= \Ric(F_{a},\cdot)^*,\end{equation} i.e.,  assuming that in local
coordinates $(x^1,\ldots,x^n)$, we have
\[ F_{a}=F^{i}_{a}\frac{\partial}{\partial x_{i}}, \]
then the evolution equation is
\[\frac{\partial F^{i}_{a}}{\partial t} = g^{ij}\Ric_{jk}F^{k}_{a}.\]
 Since this is a linear ODE, there are unique
solutions for all times $t\in [0,T)$.


The next remark to make is that this frame remains orthonormal:

\begin{claim}[Orthonormal frame remains orthonormal]\label{claim:orthonormal-frame-preservation}
\uses{def:ricci-flow}
\label{ortho}
\lean{MorganTianLib.EvolvingMetricData, MorganTianLib.EvolvingFrameData, MorganTianLib.evolvingFrame_pairing_hasDerivAt_zero, MorganTianLib.evolvingFrame_pairing_eq_initial, MorganTianLib.IsOrthonormalEvolvingFrame, MorganTianLib.evolvingFrame_isOrthonormal_of_initial}
\leanok

 Suppose that ${\mathcal F}(0)=\{F_a\}_a$ is a local  $g(0)$-orthonormal frame,
  and suppose that ${\mathcal F}(t)$ evolves according to Equation~(\ref{frameevol}). Then
 for all $t\in [0,T)$ the frame ${\mathcal F}(t)$ is a local $g(t)$-orthonormal frame.
\end{claim}

\begin{proof}
\sketch Differentiate the metric pairing using the frame evolution and
the Ricci flow equation.
\begin{eqnarray*}
\frac{\partial}{\partial t}\langle F_a(t),F_b(t)\rangle_{g(t)} & = & \langle
\frac{\partial F_a}{\partial t},F_b\rangle + \langle F_b,\frac{\partial
F_b}{\partial t}\rangle +
\frac {\partial g}{\partial t}(F_a,F_b) \\
 & = & \Ric(F_a,F_b)+\Ric(F_b,F_a)-2\Ric(F_a,F_b)=0.
\end{eqnarray*}
\end{proof}

Notice that if ${\mathcal F}'(0)=\{F'_a\}_a$ is another
frame related to ${\mathcal F}(0)$ by, say,
$$F'_a=A^b_aF_b$$
then $$F_a(t)=A^b_aF_b(t).$$ This means that the evolution of frames
actually defines a bundle automorphism
$$\Phi\colon TM|_U\times [0,T)\to TM|_U\times [0,T)$$
covering the identity map of $U\times[0,T)$ which is independent of
the choice of initial frame and is the identity at time $t=0$. Of
course, since the resulting bundle automorphism is independent of
the initial frame it globalizes to produce a bundle isomorphism
$$\Phi\colon TM\times[0,T)\to TM\times [0,T)$$
covering the identity on $M\times [0,T)$. We view this as an
evolving identification $\Phi_t$ of $TM$ with itself which is the
identity at $t=0$. The content of Claim~\ref{claim:orthonormal-frame-preservation} is:

\begin{corollary}[Pullback metric constant]\label{cor:pullback-metric-constant}
\uses{claim:orthonormal-frame-preservation}
\lean{MorganTianLib.EvolvingTransportData, MorganTianLib.evolvingTransport_pairing_hasDerivAt_zero, MorganTianLib.evolvingTransport_pairing_eq_initial, MorganTianLib.evolvingTransport_pullback_metric_constant}

$$\Phi_t^*(g(t))=g(0).$$
\end{corollary}

Returning to the local situation of the orthonormal frame ${\mathcal
F}$, we set ${\mathcal F}^*=\{F^1,\ldots, F^n\}$ equal the dual
coframe to $\{F_1,\ldots,F_n\}$. In this coframe the Riemann
curvature tensor is given by $R_{abcd}F^aF^bF^cF^d$ where
\begin{equation}\label{framecoord} R_{abcd} =
R_{ijkl}F^{i}_{a}F^{j}_{b}F^{k}_{c}F^{l}_{d}.\end{equation} One
advantage of working in the evolving frame is that the evolution
equation for the Riemann curvature tensor simplifies:


\begin{lemma}[Curvature evolution in orthonormal frame]\label{lem:curvature-evolution-frame}
\uses{def:ricci-flow}
\lean{MorganTianLib.EvolvingCurvatureData, MorganTianLib.evolvingFrame_curvatureComponent_hasDerivAt, MorganTianLib.exists_chartFrame_hasDerivAt_chartRiemannCoefOnE_sq_curvatureEvolution_of_isRicciFlowOn, MorganTianLib.evolvingFrameCurvatureRicciAction, MorganTianLib.evolvingFrameCurvatureRicciAction_apply_explicit, MorganTianLib.curvatureEvolutionCorrection_add_evolvingFrameCurvatureRicciAction, MorganTianLib.hasDerivAt_ricciFlowFrameCurvatureComponent_of_isRicciFlowOn, MorganTianLib.curvatureB_eq_sum_orthonormalBasis}
\label{M}

Suppose that the orthonormal frame ${\mathcal F}(t)$  evolves by
Formula~(\ref{frameevol}). Then we have the evolution equation
$$\frac{\partial R_{abcd}}{\partial t} = \triangle R_{abcd} +
2(B_{abcd} +B_{acbd}-B_{abdc}-B_{adbc}),$$ where $B_{abcd} =
\sum_{e,f}R_{aebf}R_{cedf}$.
\end{lemma}


\begin{proof}
\sketch The formula follows from Theorem 2.1 in \cite{Hamiltonharnack}.
\end{proof}

Of course, the other way to describe all of this  is to consider the
four-tensor $\Phi_t^*({\mathcal R}_{g(t)})=R_{abcd}F^aF^bF^cF^d$ on
$M$. Since $\Phi_t$ is a bundle map but not a bundle map induced by
a diffeomorphism, even though the pullback of the metric
$\Phi_t^*g(t)$ is constant, it is not the case that the pullback of
the curvature $\Phi_t^*{\mathcal R}_{g(t)}$ is constant.
 The next proposition gives the evolution equation for the pullback
 of the Riemann curvature tensor.



It simplifies the notation somewhat to work directly with a basis of
$\wedge^2TM$. We chose an orthonormal basis
$$\{\varphi^{1},\ldots,\varphi^{\frac{n(n-1)}{2}}\},$$
of $\wedge^2T^*_pM$ where  we have \[ \varphi^{\alpha}(F_{a},F_{b})
= \varphi^{\alpha}_{ab} \] and write the curvature tensor in this
basis as ${\mathcal T} = ({\mathcal T}_{\alpha\beta})$ so that
\begin{equation}\label{Mdefn}
R_{abcd}={\mathcal
T}_{\alpha\beta}\varphi^{\alpha}_{ab}\varphi^{\beta}_{cd}.\end{equation}


\begin{proposition}[Curvature operator evolution]\label{prop:curvature-operator-evolution}
\uses{def:curvature-operator}
\lean{MorganTianLib.curvatureOperatorSquare, MorganTianLib.curvatureOperatorSharp, MorganTianLib.curvatureOperatorSquare_isSymm_of_isSymm, MorganTianLib.curvatureOperatorSharp_isSymm_of_isSymm, MorganTianLib.curvatureOperatorReaction_isSymm_of_isSymm, MorganTianLib.curvatureOperatorEvolutionRhs_isSymm_of_isSymm, MorganTianLib.IsCurvatureOperatorEvolution}
\label{Mevol}

The evolution of the curvature operator ${\mathcal
T}(t)=\Phi_t^*\Rm(g(t))$ is given by
\[ \frac{\partial {\mathcal T}_{\alpha\beta}}{\partial t} = \triangle
{\mathcal T}_{\alpha\beta} + {\mathcal T}^{2}_{\alpha\beta} +
{\mathcal T}^{\sharp}_{\alpha\beta},
\] where ${\mathcal T}^{2}_{\alpha\beta}={\mathcal T}_{\alpha\gamma}{\mathcal
T}_{\gamma\beta}$ is the operator square; ${\mathcal T}^{\sharp}_{\alpha\beta}
 =
c_{\alpha\gamma\zeta}c_{\beta\delta\eta}{\mathcal T}_{\gamma\delta}{\mathcal
T}_{\zeta\eta}$ is the Lie algebra square; and $c_{\alpha\beta\gamma} =
\langle[\varphi^{\alpha},\varphi^{\beta}],\varphi^{\gamma}\rangle $ are the
structure constants of the Lie algebra $\mathrm{so}(n)$ relative to the basis
$\{\varphi^{\alpha}\}$. The structure constants $c_{\alpha\beta\gamma}$  are
fully antisymmetric in the three indices.
\end{proposition}

\begin{proof}
\sketch Use the Bianchi identity and expand the curvature products in the
orthonormal $\wedge^2$ basis.
\uses{lem:curvature-evolution-frame}
We work in local coordinates that are orthonormal at
the point. By the first Bianchi identity $$R_{abcd}+ R_{acdb}+
R_{adbc} = 0,$$ we get
\begin{align*}
\sum_{e,f}R_{abef}R_{cdef} &=
\sum_{e,f}(-R_{aefb}-R_{afbe})(-R_{cefd}-R_{cfde})\\& =
\sum_{e,f}2R_{aebf}R_{cedf}-2R_{aebf}R_{cfde}\\& = 2(B_{abcd} -
B_{abdc}).
\end{align*}
Note that
\begin{align*}
\sum_{e,f}R_{abef}R_{cdef}& =
\sum_{e,f}{\mathcal T}_{\alpha\beta}\varphi^{\alpha}_{ab}\varphi^{\beta}_{ef}
{\mathcal T}_{\gamma\lambda}\varphi^{\gamma}_{cd}\varphi^{\lambda}_{ef}\\
&={\mathcal T}_{\alpha\beta}\varphi^{\alpha}_{ab}{\mathcal
T}_{\gamma\lambda}\varphi^{\gamma}_{cd}\delta^{\beta\lambda}\\& =
{\mathcal
T}^{2}_{\alpha\beta}\varphi^{\alpha}_{ab}\varphi^{\beta}_{cd}.
\end{align*}
Also,
\begin{align*}
2(B_{acbd}-B_{adbc}) &= 2\sum_{e,f}(R_{aecf}R_{bedf} -
R_{aedf}R_{becf})\\& = 2\sum_{e,f}({\mathcal
T}_{\alpha\beta}\varphi^{\alpha}_{ae}\varphi^{\beta}_{cf}{\mathcal
T}_{\gamma\lambda}\varphi^{\gamma}_{be}\varphi^{\lambda}_{df} -
{\mathcal
T}_{\alpha\beta}\varphi^{\alpha}_{ae}\varphi^{\beta}_{df}{\mathcal
T}_{\gamma\lambda}\varphi^{\gamma}_{be}\varphi^{\lambda}_{cf})
\\&=
2\sum_{e,f}{\mathcal T}_{\alpha\beta}{\mathcal
T}_{\gamma\lambda}\varphi^{\alpha}_{ae}\varphi^{\gamma}_{be}(\varphi^{\beta}_{cf}\varphi^{\lambda}_{df}-\varphi^{\beta}_{df}\varphi^{\lambda}_{cf})
\\&=2\sum_e{\mathcal T}_{\alpha\beta}{\mathcal T}_{\gamma\lambda}\varphi^{\alpha}_{ae}
\varphi^{\gamma}_{be}[\varphi^{\beta},\varphi^{\lambda}]_{cd}
\\&=
\sum_e{\mathcal T}_{\alpha\beta}{\mathcal
T}_{\gamma\lambda}[\varphi^{\beta},\varphi^{\lambda}]_{cd}
(\varphi^{\alpha}_{ae}\varphi^{\gamma}_{be}
- \varphi^{\alpha}_{be}\varphi^{\gamma}_{ae})
\\&={\mathcal T}_{\alpha\beta}{\mathcal T}_{\gamma\delta}[\varphi^{\beta},\varphi^{\lambda}]_{cd}[\varphi^{\alpha},\varphi^{\gamma}]_{ab}
\\&=
{\mathcal
T}^{\sharp}_{\alpha\beta}\varphi^{\alpha}_{ab}\varphi^{\beta}_{cd}.
\end{align*}
So we can rewrite the equation for the evolution of the curvature
tensor given in Lemma~\ref{lem:curvature-evolution-frame} as \[ \frac{\partial R_{abcd}
}{\partial t} = \triangle R_{abcd} + {\mathcal
T}^{2}_{\alpha\beta}\varphi^{\alpha}_{ab}\varphi^{\beta}_{cd} +
{\mathcal
T}^{\sharp}_{\alpha\beta}\varphi^{\alpha}_{ab}\varphi^{\beta}_{cd},
\] or equivalently as \[ \frac{\partial {\mathcal T}_{\alpha\beta}}{\partial t}
= \triangle {\mathcal T}_{\alpha\beta} + {\mathcal
T}^{2}_{\alpha\beta} + {\mathcal T}^{\sharp}_{\alpha\beta}. \] We
abbreviate the last equation as
\[ \frac{\partial {\mathcal T}}{\partial t}= \triangle {\mathcal T} + {\mathcal T}^{2} +
{\mathcal T}^{\sharp}.\]
\end{proof}

\begin{remark}[Bianchi identity and operator squares]\label{rem:bianchi-identity-operator-squares}
\lean{MorganTianLib.quadraticCurvatureB, MorganTianLib.curvatureOperatorSquareTensor, MorganTianLib.curvatureOperatorSharpTensor, MorganTianLib.curvatureOperatorReactionTensor, MorganTianLib.curvatureOperatorSquareTensor_eq_two_sub, MorganTianLib.curvatureOperatorReactionTensor_bianchi, MorganTianLib.kulkarniNomizu, MorganTianLib.kulkarniNomizu_antisymm₃₄, MorganTianLib.kulkarniNomizu_pairSwap, MorganTianLib.kulkarniNomizu_bianchi, MorganTianLib.bianchiCounterexample, MorganTianLib.bianchiCounterexample_antisymm₃₄, MorganTianLib.bianchiCounterexample_pairSwap, MorganTianLib.bianchiCounterexample_bianchi, MorganTianLib.bianchiCounterexample_square_defect, MorganTianLib.bianchiCounterexample_sharp_defect}
\leanok
Notice that neither ${\mathcal T}^{2}$ nor ${\mathcal T}^{\sharp}$
satisfies the Bianchi identity, but their sum does.
\end{remark}


\section{Variation of distance under Ricci flow}

There is one result that we will use several times in the arguments
to follow. Since it is an elementary result (though the proof is
somewhat involved), we have chosen to include it here.

\begin{proposition}[Distance variation lower bound]\label{prop:distance-variation-lower-bound}
\uses{def:ricci-flow}
\lean{MorganTianLib.ForwardDiffQuotientGE.to_neg, MorganTianLib.lowerEndpointBound_of_forwardDiffQuotientGE, MorganTianLib.fixedCurveLength_forwardDiffQuotientGE_of_hasDerivWithinAt, MorganTianLib.metricInner_forwardDiffQuotientGE_of_metricVariation, MorganTianLib.metricInner_forwardDiffQuotientGE_of_isRicciFlowOn}
\label{I.8.3}

 Let $t_0\in \Ar$ and let $(M,g(t))$ be a Ricci flow defined for $t$ in an interval containing $t_0$
 with $(M,g(t))$ complete for every $t$ in this interval. Fix a constant $K<\infty$.  Let $x_0,x_1$ be two points
of $M$ and let $r_0>0$ such that $d_{t_0}(x_0,x_1)\ge 2r_0$. Suppose
that $\Ric(x,t_0)\le (n-1)K$ for all $x\in B(x_0,r_0,t_0)\cup
B(x_1,r_0,t_0)$. Then
$$\frac{d(d_t(x_0,x_1))}{dt}\Bigl|_{t=t_0}\Bigr.\ge -2(n-1)\left(\frac{2}{3}Kr_0+r_0^{-1}\right).$$
 If the distance function $d_{t}(x_0,x_1)$ is not a differentiable
function of $t$ at $t=t_0$ then this inequality is understood as an
inequality for the forward difference quotient.
\end{proposition}


\begin{remark}[Forward difference quotient]\label{rem:forward-difference-quotient}
\lean{MorganTianLib.ForwardDiffQuotientGE, MorganTianLib.HasDerivWithinAt.forwardDiffQuotientGE, MorganTianLib.HasDerivAt.le_deriv_iff_forwardDiffQuotientGE}
\leanok
Of course, if the distance function is differentiable at $t=t_0$
then the derivative statement is equivalent to the forward
difference quotient statement. Thus, in the proof of this result we
shall always work with the forward difference quotients.
\end{remark}

\begin{proof}
  \uses{def:geodesic}
\sketch Apply Claim~\ref{claim:geodesic-length-suffices} to the estimate in
Claim~\ref{claim:minimal-geodesic-time-derivative}.

\end{proof}

\begin{claim}[Geodesic length suffices for distance]\label{claim:geodesic-length-suffices}
\uses{def:geodesic}
\label{geosuff}

Suppose that for every minimal $g(t_0)$-geodesic $\gamma$ from $x_0$
to $x_1$ the function $\ell_t(\gamma)$ which is the $g(t)$-length of
$\gamma$ satisfies
$$\frac{d(\ell_t(\gamma))}{dt}\Bigl|_{t=t_0}\Bigr.\ge C.$$
Then
$$\frac{d(d_t(x_0,x_1))}{dt}\Bigl|_{t=t_0}\Bigr.\ge C,$$
where, as in the statement of the proposition, if the distance
function is not differentiable at $t_0$ then the inequality in the
conclusion is interpreted by replacing the derivative on the
left-hand side with the $\liminf$ of the forward difference
quotients of $d_t(x_0,x_1)$ at  $t_0$.
\end{claim}

\begin{proof}
\sketch Approximate the distance at nearby times by minimal geodesics at
$t_0$ and take the lower limit of the forward difference quotients.
\end{proof}



The second step in the proof is to estimate the time derivative of a
minimal geodesic under the hypothesis of the proposition.

\begin{claim}[Minimal geodesic time derivative]\label{claim:minimal-geodesic-time-derivative}
\uses{def:geodesic}
\label{3.23}

Assuming the hypothesis of the proposition, for any
 minimal $g(t_0)$-geodesic $\gamma$ from $x_0$ to $x_1$, we have
$$\frac{d(\ell_t(\gamma))}{dt}\Bigl|_{t=t_0}\Bigr.\ge -2(n-1)
\left(\frac{2}{3}Kr_0+r_0^{-1}\right).$$
\end{claim}

\begin{proof}
\sketch Fix a minimal $g(t_0)$-geodesic $\gamma(u)$ from $x_0$ to $x_1$,
parameterized by arc length. We set $d=d_{t_0}(x_0,x_1)$, we set
$X(u)=\gamma'(u)$, and we take tangent vectors $Y_1,\ldots,Y_{n-1}$
in $T_{x_0}M$ which together with $X(0)=\gamma'(0)$ form an
orthonormal basis. We let $Y_i(u)$ be the parallel translation of
$Y_i$ along $\gamma$. Define $f\colon [0,d]\to [0,1]$ by:
$$f(u)=\begin{cases}u/r_0  & 0\le u\le r_0 \\ 1 & r_0\le u\le d-r_0 \\
(d-u)/r_0 & d-u\le r_0\le d, \end{cases}$$ and define $$\widetilde
Y_i(u)=f(u)Y_i(u).$$ See the figure in [Morgan--Tian].
 For $1\le i\le n-1$, let $s''_{\widetilde
Y_i}(\gamma)$ be the second variation of the $g(t_0)$-length of
$\gamma$ along $\widetilde Y_i$. Since $\gamma$ is a minimal
$g(t_0)$-geodesic, for all $i$ we have
\begin{equation}\label{geoineq}
s''_{\widetilde Y_i}(\gamma)\ge 0.\end{equation}




Let us now compute $s''_{\widetilde Y_i}(\gamma)$ by taking a
two-parameter family $\gamma(u,s)$ such that the curve $\gamma(u,0)$
is the original minimal geodesic and $\frac{\partial}{\partial
s}(\gamma(u,s))|_{s=0}=\widetilde Y_i(u)$.  We denote by $X(u,s)$
the image $D\gamma_{(u,s)}(\partial/\partial u)$ and by $\widetilde
Y_i(u,s)$ the image $D\gamma_{(u,s)}(\partial/\partial s)$. We wish
to compute
\begin{eqnarray}
\lefteqn{s''_{\widetilde Y_i}(\gamma) =
\frac{d^2}{ds^2}\left(\int_0^d\sqrt{X(u,s),X(u,s)}du\right)\Bigl|_{s=0}\Bigr.}
& &
 \nonumber \\
& = & \frac{d}{ds}\left(\int_0^d\langle X(u,s),X(u,s)\rangle^{-1/2}
\langle
X(u,s),\nabla_{\widetilde Y_i}X(u,s)\rangle du\right)\Bigl|_{s=0}\Bigr. \nonumber \\
& = &  \int_0^d-\langle X(u,0),X(u,0)\rangle^{-3/2}\langle
X(u,0),\nabla_{\widetilde
Y_i}X(u,0)\rangle^2du \label{eq2ndvar} \\
& &  +\int_0^d\frac{\langle \nabla_{\widetilde
Y_i}X(u,0),\nabla_{\widetilde Y_i}X(u,0)\rangle +\langle
X(u,0),\nabla_{\widetilde Y_i}\nabla_{\widetilde
Y_i}X(u,0)\rangle}{\langle X(u,0),X(u,0)\rangle^{1/2}}du \nonumber.
\end{eqnarray}
Using the fact that $X$ and $\widetilde Y_i$ commute (since they are
the coordinate partial derivatives of a map of a surface into $M$)
and using the fact that $Y_i(u)$ is parallel along $\gamma$, meaning
that $\nabla_X(Y_i)(u)=0$, we see that $\nabla_{\widetilde
Y_i}X(u,0)=\nabla_X\widetilde Y_i(u,0)=f'(u)Y_i(u)$. By construction
$\langle Y_i(u),X(u,0)\rangle=0$. It follows that
$$\langle \nabla_{\widetilde Y_i}X(u,0),X(u,0)\rangle=\langle \nabla_X(\widetilde
Y_i)(u,0),X(u,0)\rangle = \langle f'(u)Y_i(u),X(u,0)\rangle =0.$$
Also, $\langle X(u,0),X(u,0)\rangle=1$, and by construction $\langle
Y_i(u,0),Y_i(u,0)\rangle=1$. Thus, Equation~(\ref{eq2ndvar})
simplifies to
\begin{eqnarray}
\lefteqn{s''_{\widetilde
Y_i}(\gamma)=\frac{d^2}{ds^2}\left(\int_0^d\sqrt{X(u,s),X(u,s)}du\right)\Bigl|_{s=0}\Bigr.
 \label{2ndvaria}}
\\& = & \int_0^d\left((f'(u))^2\langle
Y_i(u),Y_i(u)\rangle+\langle \nabla_{\widetilde
Y_i}\nabla_X(\widetilde
Y_i(u,0)),X(u,0)\rangle\right)du \nonumber\\
& = & \int_0^d\left(\langle R(\widetilde Y_i,X)\widetilde
Y_i(u,0),X(u,0)\rangle-\langle \nabla_X\nabla_{\widetilde
Y_i}\widetilde Y_i(u,0),X(u,0)\rangle+(f'(u))^2\right)
du.\nonumber\end{eqnarray} Now we restrict to $s=0$ and for simplicity of notation
we leave the variable $u$ implicit. We have
$$\langle
\nabla_X\nabla_{\widetilde Y_i}\widetilde
Y_i,X\rangle=\frac{d}{du}\langle \nabla_{\widetilde Y_i}\widetilde
Y_i,X\rangle-\langle \nabla_{\widetilde Y_i}\widetilde
Y_i,\nabla_XX\rangle=\frac{d}{du}\langle \nabla_{\widetilde
Y_i}\widetilde Y_i,X\rangle,$$ where the last equality is a
consequence of the geodesic equation, $\nabla_XX=0$. It follows that
$$\int_0^d\langle
\nabla_X\nabla_{\widetilde Y_i}\widetilde Y_i,X\rangle
du=\int_0^d\frac{d}{du}\langle \nabla_{\widetilde Y_i}\widetilde
Y_i,X\rangle=0,$$ where the last equality is a consequence of the
fact that $\widetilde Y_i$ vanishes at the end points.

Consequently, plugging these into Equation~(\ref{2ndvaria}) we have
\begin{equation} \label{sformula} s''_{\widetilde Y_i}(\gamma) =
  \int_0^d\left(\langle R(\widetilde Y_i,X)\widetilde
Y_i(u,0),X(u,0)\rangle+(f'(u))^2\right)du. \end{equation} Of course,
it is immediate from the definition that $f'(u)^2=1/r_0^2$ for $u\in
[0,r_0]$ and for $u\in [d-r_0,d]$ and is zero otherwise. Also, from
the definition of the vector fields $Y_i$ we have
$$\sum_{i=1}^{n-1}\langle
R(Y_i,X)Y_i(u),X(u)\rangle=-\Ric_{g(t_0)}(X(u),X(u)),$$ so that
$$\sum_{i=1}^{n-1}\langle
R(\widetilde Y_i,X)\widetilde
Y_i(u),X(u)\rangle=-f^2(u)\Ric_{g(t_0)}(X(u),X(u)).$$

Hence, summing Equalities~(\ref{sformula}) for $i=1,\ldots, n-1$ and
using Equation~(\ref{geoineq}) gives
\begin{eqnarray*}
0\le\sum_{i=1}^{n-1}s''_{\widetilde Y_i}(\gamma) & = &
\int_0^{r_0}\left[\frac{u^2}{r_0^2}\left(-\Ric_{g(t_0)}\left(X(u),X(u)\right)\right)
+\frac{n-1}{r_0^2}\right]du
\\
& & + \int_{r_0}^{d-r_0}-\Ric_{g(t_0)}(X(u),X(u))du
\\
 & & +\int_{d-r_0}^d\left[\frac{(d-u)^2}{r_0^2}
\left(-\Ric_{g(t_0)}(X(u),X(u))\right)+\frac{n-1}{r_0^2}\right]du.\end{eqnarray*}
Rearranging the terms yields
\begin{eqnarray*} 0 & \le & -\int_0^d \Ric_{g(t_0)}(X(u),X(u))du \\
& &+\int_0^{r_0}\left[\left(1-\frac{u^2}{r_0^2}\right)
\left(\Ric_{g(t_0)}(X(u),X(u))\right) +\frac{n-1}{r_0^2}\right]du\\
& & +\int_{d-r_0}^d
\left[\left(1-\frac{(d-u)^2}{r_0^2}\right)\left(\Ric_{g(t_0)}(X(u),X(u))\right)
+\frac{n-1}{r_0^2}\right]du.
\end{eqnarray*}
Since
\begin{eqnarray*}
\frac{d(\ell_t(\gamma))}{dt}\Bigl|_{t=t_0}\Bigr. & = &
\frac{d}{dt}\left[\left(\int_0^d\sqrt{\langle X(u),X(u)\rangle}dt\right)^{1/2}\right]|_{t=t_0} \\
& = & -\int_0^d \Ric_{g(t_0)}(X(u),X(u))du,
\end{eqnarray*}
we have
\begin{eqnarray*}
\lefteqn{\frac{d(\ell_t(\gamma))}{dt}\Bigl|_{t=t_0}\Bigr.\ge
-\left\{\int_0^{r_0}\left[\left(1-\frac{u^2}{r_0^2}\right)\left(\Ric_{g(t_0)}(X(u),X(u))\right)
+\frac{n-1}{r_0^2}\right]du\right.}
\\
& & \left. +\int_{d-r_0}^d
\left[\left(1-\frac{(d-u)^2}{r_0^2}\right)\left(\Ric_{g(t_0)}(X(u),X(u))\right)
+\frac{n-1}{r_0^2}\right]du\right\}
\end{eqnarray*}


Now, since $|X(u)|=1$, by the hypothesis of the proposition we have
the estimate $\Ric_{g(t_0)}(X(u),X(u))\le (n-1)K$ on the regions of
integration on the right-hand side of the above inequality. Thus,
\begin{eqnarray*}
\frac{d(\ell_t(\gamma))}{dt}\Bigl|_{t=t_0}\Bigr.\ge
-2(n-1)\left(\frac{2}{3}r_0K+r_0^{-1}\right).
\end{eqnarray*}
This completes the proof of Claim~\ref{claim:minimal-geodesic-time-derivative}.
\end{proof}

\begin{corollary}[Distance variation on complete manifolds]\label{cor:distance-variation-complete}
\uses{def:ricci-flow}
\label{2ndI.8.3}

Let $t_0\in \Ar$ and let $(M,g(t))$ be a Ricci flow defined for $t$ in an
interval
containing $t_0$ and with $(M,g(t))$ complete for every $t$ in this
interval. Fix a constant $K<\infty$. Suppose that $\Ric(x,t_0)\le (n-1)K$
for all $x\in M$.
 Then for any points $x_0,x_1\in M$ we have
$$\frac{d(d_t(x_0,x_1))}{dt}\Bigl|_{t=t_0}\Bigr.\ge -4(n-1)\sqrt{\frac{2K}{3}}$$
in the sense of forward difference quotients.
\end{corollary}

\begin{proof}
\sketch Use the two distance cases and apply the preceding variation bounds.
\uses{prop:distance-variation-lower-bound,claim:geodesic-length-suffices}
There are two cases: Case (i): $d_{t_0}(x_0,x_1)\ge
\sqrt{\frac{6}{K}}$ and Case (ii) $d_{t_0}(x_0,x_1)<
\sqrt{\frac{6}{K}}$. In Case (i) we  take $r_0=\sqrt{3/2K}$ in
Proposition~\ref{prop:distance-variation-lower-bound}, and we conclude that the $\liminf$ at
$t_0$ of the difference quotients for $d_{t}(x_0,x_1)$ is at most
$-4(n-1)\sqrt{\frac{2K}{3}}$. In Case (ii) w let $\gamma(u)$ be any
minimal $g(t_0)$-geodesic from $x_0$ to $x_1$ parameterized by arc
length. Since
$$\frac{d}{dt}(\ell_t(\gamma))|_{t=t_0}=-\int_\gamma
\Ric_{g(t_0)}(\gamma'(u),\gamma'(u))du,$$ we see that
$$\frac{d}{dt}(\ell_t(\gamma))|_{t=t_0}\ge
-(n-1)K\sqrt{6/K}=-(n-1)\sqrt{6K}.$$ By Claim~\ref{claim:geodesic-length-suffices}, this
implies that the $\liminf$ of the forward difference quotient of
$d_t(x_0,x_1)$ at $t=t_0$ is at least $-(n-1)\sqrt{6K}\ge
-4(n-1)\sqrt{2K/3}$.
\end{proof}



\begin{corollary}[Distance integral bound]\label{cor:distance-integral-bound}
\uses{def:ricci-flow}
\label{corI.8.3}

 Let $(M,g(t)),\ a\le t\le b$, be a Ricci flow with $(M,g(t))$ complete for every $t
\in [0,T)$. Fix a positive function $K(t)$, and suppose that $\Ric_{g(t)}(x,t)\le (n-1)K(t)$ for all $x\in M$ and all $t\in [a,b]$. Let
$x_0,x_1$ be two points of $M$. Then
$$d_a(x_0,x_1)\le d_b(x_0,x_1)+4(n-1)\int_a^b\sqrt{\frac{2K(t)}{3}}dt.$$
\end{corollary}


\begin{proof}
\sketch Integrate the forward difference-quotient inequality in time.
\uses{cor:distance-variation-complete,lem:forward-difference-quotient}
By Corollary~\ref{cor:distance-variation-complete} we have
\begin{equation}\label{liminf}
\frac{d}{dt}d_t(x_0,x_1)|_{t=t'}\ge -4(n-1)\sqrt{\frac{2K(t')}{3}}
\end{equation}
in the sense  of forward difference quotients. Thus, this result is
an immediate consequence of Lemma~\ref{lem:forward-difference-quotient}.
\end{proof}


\section{Shi's derivative estimates}

The last `elementary' result we discuss is Shi's result controlling all derivatives in terms of a
bound on curvature. This is a consequence of the parabolic nature of
the Ricci flow equation. More precisely, we can control all
derivatives of the curvature tensor at a point $p\in M$ and at a
time $t$ provided that we have an upper bound for the curvature on
an entire backward parabolic neighborhood of $(p,t)$ in space-time.
The estimates become weaker as the parabolic neighborhood shrinks,
either in the space direction or the time direction.

Recall that for any $K<\infty$  if $(M,g)$ is a Riemannian manifold
with $|\Rm|\le K$ and if for some $r\le \pi/\sqrt{K}$ the metric ball
$B(p,r)$ has compact closure in $M$, then the exponential mapping
$\exp_p$ is defined on the ball $B(0,r)$ of radius $r$ centered
at the origin of $T_pM$ and $\exp_p\colon B(0,r)\to M$ is a
local diffeomorphism onto $B(p,r)$.


The first of Shi's derivative estimates controls the first
derivative of $\Rm$.

\begin{theorem}[Shi's first derivative estimate]\label{thm:shi-first-derivative}
\uses{def:ricci-flow}
\lean{MorganTianLib.shiCoefficient, MorganTianLib.shiCoefficient_pos, MorganTianLib.shiCoefficient_telescope, MorganTianLib.shiTowerCombination, MorganTianLib.shiTopTerm_le_tower, MorganTianLib.shiReactionSum_le, MorganTianLib.ShiTowerInequalitiesOn, MorganTianLib.shiTowerCombination_deriv_le, MorganTianLib.shiTower_sqrt_le_div_on_compact_interior, MorganTianLib.exists_uniform_shiTower_sqrt_bound_on_compact_interior, MorganTianLib.riemannCovDerivTower, MorganTianLib.riemannCovDerivNormSqAt, MorganTianLib.riemannCovDerivNormAt, MorganTianLib.riemannCovDerivNormSqAt_nonneg, MorganTianLib.riemannCovDerivNormAt_sq, MorganTianLib.covTensorNormSqAt_covTensorDeriv_eq_sum, MorganTianLib.riemannCovDerivNormSqAt_succ_eq_sum, MorganTianLib.covTensorNormSqAt_covTensorDerivAlong_le, MorganTianLib.covTensorNormSqAt_covTensorDeriv_eq_zero_iff_directional, MorganTianLib.riemannCovDerivNormSqAt_succ_directional_le, MorganTianLib.riemannCovDerivNormSqAt_zero_eq_curvatureForm_sq_sum, MorganTianLib.curvatureForm_sq_sum_nonneg}
\label{firstshi}

There is a constant $C=C(n)$, depending only on the dimension $n$,
such that the following holds for every $K<\infty$, for every $T>0$,
and for every $r>0$. Suppose that $(U,g(t)),\ 0\le t\le T$, is an $n$-dimensional
Ricci flow with $|\Rm(x,t)|\le K$ for all $x\in U$ and $t\in [0,T]$.
Suppose that $p\in U$ has the property that $B(p,0,r)$ has compact
closure in $U$. Then
$$|\nabla \Rm(p,t)|\le
CK\left(\frac{1}{r^2}+\frac{1}{t}+K\right)^{1/2}.$$
\end{theorem}


For a proof of this result see Chapter 6.2, starting on page 212, of
\cite{ChowLuNi}.

We also need higher derivative estimates. These are also due to Shi,
but they take a slightly different form; see \cite{Shi1,Shi2}. (See Theorem 6.9 on page
210 of \cite{ChowLuNi}.)


\begin{theorem}[Shi's derivative estimates]\label{thm:shi-derivative-estimates}
\uses{def:ricci-flow}
\lean{MorganTianLib.ShiTowerInequalitiesOn, MorganTianLib.shiTowerCombination_hasDerivWithinAt, MorganTianLib.shiTowerCombination_deriv_le, MorganTianLib.shiTower_mul_pow_le_affine, MorganTianLib.shiTower_sqrt_le_div, MorganTianLib.shiTower_sqrt_le_div_on_compact, MorganTianLib.shiTower_sqrt_le_div_on_compact_interior, MorganTianLib.exists_uniform_shiTower_sqrt_bound_on_compact_interior, MorganTianLib.riemannCovDerivTower, MorganTianLib.riemannCovDerivNormSqAt, MorganTianLib.riemannCovDerivNormAt, MorganTianLib.riemannCovDerivNormSqAt_nonneg, MorganTianLib.riemannCovDerivNormAt_sq, MorganTianLib.covTensorNormSqAt_covTensorDeriv_eq_sum, MorganTianLib.riemannCovDerivNormSqAt_succ_eq_sum, MorganTianLib.covTensorNormSqAt_covTensorDerivAlong_le, MorganTianLib.covTensorNormSqAt_covTensorDeriv_eq_zero_iff_directional, MorganTianLib.riemannCovDerivNormSqAt_succ_directional_le, MorganTianLib.shiEuclideanCutoff_exists_derivative_bounds, MorganTianLib.abs_covTensorDerivAlongComponentAt_le_shi_bound_on_compact_interior}
\label{shi}

{\bf (Shi's Derivative Estimates)} Fix the dimension $n$ of the Ricci flows under consideration. Let
$K<\infty$ and $\alpha>0$ be positive constants.  Then for each
non-negative integer $k$ and each $r>0$ there is a constant
$C_k=C_k(K,\alpha,r,n)$ such that the following holds. Let
$(U,g(t)),\ 0\le t\le T$, be a Ricci flow with $T\le \alpha/K$.
Fix $p\in U$ and suppose that the metric ball $B(p,0,r)$ has compact
closure in $U$. If
$$|\Rm(x,t)|\le K\ \ \text{for all }\ (x,t)\in P(x,0,r,T),$$
then
$$|\nabla^k(\Rm(y,t))|\le \frac{C_k}{t^{k/2}}$$
for all $y\in B(p,0,r/2)$ and all $t\in(0,T]$.
\end{theorem}

For a proof of this result see Chapter 6.2 of \cite{ChowLuNi} where
these estimates are proved for the first and second derivatives of
$\Rm$. The proofs of the higher derivatives follow similarly.  Below, we
shall prove a stronger form of this result below including the proof
for all derivatives.

We shall need a stronger version of this result, a version which is
well-known but for which there seems to be no good reference. The
stronger version takes as hypothesis $C^k$-bounds on the
initial conditions and produces a better bound on the derivatives of
the curvature at later times. The argument is basically the same as
that of the result cited above, but since there is no good reference
for it we include the proof, which was shown to us by Lu Peng.





\begin{theorem}[Shi's derivative estimates with initial bounds]\label{thm:shi-with-initial-derivatives}
\uses{def:ricci-flow}
\lean{MorganTianLib.shiFm_zero_le_of_level_bounds, MorganTianLib.shiFm_zero_of_le, MorganTianLib.shiFm_level_le_of_bound, MorganTianLib.shiFm_sqrt_level_le_of_bound}
\label{shiw/deriv}

Fix the dimension $n$ of the Ricci flows under consideration. Let
$K<\infty$ and $\alpha>0$ be given positive constants. Fix an
integer $l\ge 0$. Then for each integer $k\ge 0$ and for each $r>0$
there is a constant $C'_{k,l}=C'_{k,l}(K,\alpha,r,n)$ such that the
following holds. Let $(U,g(t)),\ 0\le t\le T$, be a Ricci flow with
$T\le \alpha/K$. Fix $p\in U$ and suppose that the metric ball
$B(p,0,r)$ has compact closure in $U$. Suppose that
\begin{eqnarray*}
\left\vert \Rm\left(  x,t\right)  \right\vert  &
\leq K\ \ \text{ for all }x\in U\text{ and all}\  t\in[0,T],\\
\left\vert \nabla^{\beta}\Rm\left(  x,0\right)
\right\vert  & \leq K\ \ \text{ for all }x\in U\ \ \text{ and all }\
\ \beta\leq l.
\end{eqnarray*}
Then
\[
\left\vert \nabla^{k}\Rm\left(  y,t\right)
\right\vert
\leq\frac{C'_{k,l}}{t^{\max\left\{  k-l,0\right\}  /2}}%
\]
for all $y\in$ $B(p,0,r/2)  $ and all $t\in(0,T].$ In particular if
$k\leq l$, then for $y\in B(p,0,r/2)$ and $t\in (0,T]$ we have
\[
\left\vert \nabla^{k}\Rm\left(  y,t\right)
\right\vert \leq C_{k,l}'.
\]
\end{theorem}


\begin{proof}
\sketch Induct on $k$, use the curvature derivative evolution inequality,
localize with Proposition~\ref{prop:parabolic-cutoff-function}, and apply
the maximum principle.
\uses{prop:parabolic-cutoff-function}
The first remark is that establishing Theorem~\ref{thm:shi-with-initial-derivatives} for
one value of $r$ immediately gives it for all $r'\ge 2r$. The reason
is that for such $r'$ any point  $y\in B(p,0,r'/2)$ has the property
that $B(y,0,r)\subset B(p,0,r')$ so that a curvature bound on
$B(p,0,r')$ will imply one on $B(y,0,r)$ and hence by the result for
$r$ will imply the higher derivative bounds at $y$.


 Thus, without loss of
generality we can suppose that $r\le \pi/2\sqrt{K}$. We shall assume this
from now on in the proof. Since
$B(p,0,r)$ has compact closure in $M$, for some $r<r'<\pi/\sqrt{K}$
the ball $B(p,0,r')$ also has compact closure in $M$. This means
that the exponential mapping from the ball of radius $r'$ in $T_pM$
is a local diffeomorphism onto $B(p,0,r')$.

The proof is by induction: We assume that we have established the
result for $k=0,\ldots,m$, and then we shall establish it for
$k=m+1$. The inductive hypothesis tells us that there are constants
$A_j,\ 0\le j\le m$, depending on $(l,K,\alpha,r,n)$ such that for
all $(x,t)\in B(p,0,r/2)\times (0,T]$ we have
\begin{equation}\label{strongineq}
\left\vert \nabla^{j}\Rm \left(  x,t\right)
\right\vert \leq A_{j}t^{-\max\left\{  j-l,0\right\}  /2}.
\end{equation}
 Even better, applying the inductive result to $B(y,0,r/2)$ with
$y\in B(p,0,r/2)$ we see, after replacing the $A_j$ by the
larger constants associated with $(l,K,\alpha,r/2,n)$, that we have
the same inequality for all $y\in B(x,0,3r/4)$.


We fix a constant $C\geq \max(4A^2_m,1)$ and consider
\[
F_{m}(x,t)=\left(  C+t^{\max\left\{ m-l,0\right\} }  \left\vert \nabla
^{m}\Rm(x,t)\right\vert ^{2}\right)  t^{\max\left\{
m+1-l,0\right\}  }\left\vert
\nabla^{m+1}\Rm(x,t)\right\vert ^{2}.
\]
Notice that bounding $F_m$ above by a constant $(C'_{m+1,l})^2$ will
yield
$$|\nabla^{m+1}\Rm(x,t)|^2\le \frac{(C'_{m+1,l})^2}{t^{\max\{m+1-l,0\}}},$$
and hence will complete the proof of the result.

 Bounding $F_m$
above (assuming the inductive hypothesis) is what is accomplished in
the rest of this proof. The main calculation is the proof of the
following claim under the inductive hypothesis.
\end{proof}

\begin{claim}[Function evolution bound]\label{claim:function-evolution-bound}
\lean{MorganTianLib.shiFm, MorganTianLib.shiFm_nonneg, MorganTianLib.shiFm_top_term_le, MorganTianLib.shiCutoffMaximumBound, MorganTianLib.shiCutoffMaximumBound_strong, MorganTianLib.shiCutoffMaximumBound_strong_or_small, MorganTianLib.evolvingFrameCurvatureEnergy_hasDerivAt, MorganTianLib.evolvingFrameCurvatureEnergy_abs_deriv_le_young, MorganTianLib.abs_curvatureOperatorEnergy_deriv_sub_laplacianPairing_le, MorganTianLib.shiReactionTerm_le_quadratic, MorganTianLib.shiReactionSum_le_quadratic, MorganTianLib.shiReactionSum_le_uniform, MorganTianLib.shiReactionSum_le_uniform_of_le, MorganTianLib.riemannCovDerivTower, MorganTianLib.riemannCovDerivNormSqAt, MorganTianLib.riemannCovDerivNormAt}
\label{Fm}
 With $F_m$ as defined above and with $C\ge \max(4A^2_m,1)$,
there are constants $c_1$ and $C_0,C_1$ depending on $C$ as well as
$K,\alpha,A_1,\ldots,A_m$ for which the following holds on
$B(p,0,3r/4)\times (0,T]$:
$$
\left(  \frac{\partial}{\partial t}-\Delta\right)
F_{m}(x,t)\leq-\frac {c_1}{t^{\mathrm{s}\left\{  \max\left\{
m-l+1,0\right\}  \right\} }}\left( F_{m}(x,t)-C_0\right)
^{2}+\frac{C_1}{t^{\mathrm{s}\{\max\left\{ m-l+1,0\right\} \} }},
$$
where
$$\mathrm{s}(n)=\begin{cases} +1 & \text{if } n>0 \\ 0 & \text{if }n=0 \\ -1&
\text{if }n<0.\end{cases}$$
\end{claim}

\begin{proof}
\sketch The evolution inequality follows by differentiating $F_m$ and
absorbing lower-order terms using the choice of $C$.
\end{proof}

\begin{proof}
\sketch Let us assume this claim and use it to prove Theorem~\ref{thm:shi-with-initial-derivatives}. We fix
$C=\max\{4A_m^2,1\}$, and consider the resulting function $F_m$. The
constants $c_1,C_0,C_1$ from Claim~\ref{claim:function-evolution-bound} depend only on $K,\alpha$, and
$A_1,\ldots,A_m$. Since $r\le \pi/2\sqrt{K}$, and $B(p,0,r)$ has
compact closure in $U$, there is some $r'>r$ so that the exponential
mapping $\exp_p\colon B(0,r')\to U$ is a local diffeomorphism
onto $B(p,0,r')$. Pulling back by the exponential map, we  replace
the Ricci flow on $U$ by a Ricci flow on $B(0,r')$ in $T_pM$.
Clearly, it suffices to establish the estimates in the statement of
the proposition for $B(0,r/2)$. This remark allows us to assume that
the exponential mapping is a diffeomorphism onto $B(p,0,r)$. Bounded
curvature then comes into play in the following crucial proposition,
which goes back to Shi. The function given in the next proposition allows us to localize
the computation in the ball $B(p,0,r)$.
\end{proof}

\begin{proposition}[Parabolic cutoff function]\label{prop:parabolic-cutoff-function}
\uses{def:ricci-flow}
\lean{MorganTianLib.ShiParabolicCutoff, MorganTianLib.shiEuclideanCutoff, MorganTianLib.shiEuclideanCutoff_contDiff, MorganTianLib.shiEuclideanCutoff_nonneg, MorganTianLib.shiEuclideanCutoff_le_one, MorganTianLib.shiEuclideanCutoff_eq_one_of_mem, MorganTianLib.shiEuclideanCutoff_eq_zero_of_not_mem, MorganTianLib.shiEuclideanCutoff_hasCompactSupport, MorganTianLib.shiParabolicCutoff_of_euclidean_bump, MorganTianLib.shiEuclideanCutoff_exists_derivative_bounds, MorganTianLib.shiCutoffMaximumBound, MorganTianLib.shiCutoffMaximumBound_strong, MorganTianLib.shiCutoffMaximumBound_strong_or_small, MorganTianLib.shi_le_affineBarrier_of_parabolic_inequality, MorganTianLib.shi_le_affineBarrier_on_compact}
\label{etaprop}

Fix constants $0<\alpha$ and the dimension $n$. Then there is a constant
$C_2'=C_2'(\alpha,n)$ and for each $r>0$ and $K<\infty$ there is a constant
$C_2=C_2(K,\alpha,r,n)$ such that the following holds. Suppose that $(U,g(t)),\
0\le t\le T$,  is an $n$-dimensional Ricci flow  with $T\le \alpha/K$. Suppose
that $p\in U$ and that $B(p,0,r)$ has compact closure in $U$ and that the
exponential mapping from the ball of radius $r$ in $T_pU$ to $B(p,0,r)$ is a
diffeomorphism. Suppose that $|\Rm(x,0)|\le K$ for all $x\in B(p,0,r)$. There is
a smooth function $\eta\colon B(p,0,r)\to [0,1]$ satisfying the following for
all $t\in [0,T]$:
\begin{enumerate}
\item $\eta$ has compact support in $B(p,0,r/2)$
\item The restriction of $\eta$ to $B(p,0,r/4)$ is identically $1$.
\item $|\Delta_{g(t)} \eta|\le C_2(K,\alpha,r,n)$.
\item $\frac{|\nabla \eta|_{g(t)}^2}{\eta}\le \frac{C'_2(\alpha,n)}{r^2}$
\end{enumerate}
\end{proposition}

\begin{proof}
\sketch
For a proof of this result see Lemma 6.62 on page 225 of
\cite{ChowLuNi}.

We can apply this proposition to our situation, because we are
assuming that $r\le \pi/2\sqrt{K}$ so that the exponential mapping
is a local diffeomorphism onto $B(p,0,r)$ and we have pulled the
Ricci flow back to the ball in the tangent space.

Fix any $y\in B(p,0,r/2)$ and choose $\eta$ as in the previous
proposition for the constants $C_2(\alpha,n)$ and
$C_2'(K,\alpha,r/4,n)$. Notice that $B(y,0,r/4)\subset B(p,0,3r/4)$
so that the conclusion of Claim~\ref{claim:function-evolution-bound} holds for every $(z,t)$
with $z\in B(y,0,r/4)$ and $t\in [0,T]$.  We shall show that the restriction of $\eta
F_m$ to $P(y,0,r/4,T)$ is bounded by a constant that depends only on
$K,\alpha,r,n,A_1,\ldots,A_m$. It will then follow immediately that
the restriction of $F_m$ to $P(y,0,r/8,T)$ is bounded by the same
constant. In particular, the values of $F_m(y,t)$ are bounded by the
same constant for all $y\in B(p,0,r/2)$ and $t\in [0,T]$.

 Consider a point $(x,t)\in B(y,0,r/2)\times
[0,T]$ where $\eta F_m$ achieves its maximum; such a point exists
since the ball $B(y,0,r/2)\subset B(p,0,r)$, and hence $B(y,0,r/2)$
has compact closure in $U$. If $t=0$, then $\eta F_m$ is bounded by
$(C+K^2)K^2$ which is a constant depending only on $K$ and $A_m$.
This, of course, immediately implies the result. Thus we can assume
that the maximum is achieved at some $t>0$. When $s\left\{
\max\left\{ m+1-l,0\right\} \right\}  -0,$ according to the
Claim~\ref{claim:function-evolution-bound}, we have

\[
\left(  \frac{\partial}{\partial t}-\Delta\right)
F_{m}\leq-c_1\left( F_{m}-C_0\right)  ^{2}+C_1.
\]

We compute
\[
\left(  \frac{\partial}{\partial t}-\Delta\right)  \left(  \eta
F_{m}\right) \leq\eta\left(  -c_1\left(
F_{m-C_0}\right)^{2}+C_1\right) -\Delta\eta\cdot
F_m-2\nabla\eta\cdot\nabla F_m.
\]


 Since $(x,t)$ is a maximum point for $\eta F_m$ and since
$t>0$,  a simple maximum principle argument shows that
$$\left(  \frac{\partial}{\partial t}-\Delta\right)\eta  F_{m}(x,t)\ge 0.$$
Hence, in this case we conclude that
\begin{eqnarray*} 0& \le & \left(
\frac{\partial}{\partial t}-\Delta\right)  \left( \eta(x)
F_{m}(x,t)\right) \leq\eta(x)\left(  -c_1\left(
F_{m}(x,t)-C_0\right)^{2}+C_1\right) \\
& &  -\Delta\eta(x)\cdot F_m(x,t)-2\nabla\eta(x)\cdot\nabla
F_m(x,t). \end{eqnarray*}
 Hence, $$c_1\eta (F_m(x,t)-C_0)^2\le \eta(x)
C_1-\Delta\eta(x)\cdot F_m(x,t)-2\nabla\eta(x)\cdot \nabla F_m(x,t).$$ Since we
are proving that $F_m$ is bounded, we are free to argue by contradiction and
assume that $F_m(x,t)\ge 2C_0$, implying that $F_m(x,t)-C_0\ge F_m(x,t)/2$.
Using this inequality yields \begin{eqnarray*} \eta(x) (F_m(x,t)-C_0) & \le &
\frac{2\eta C_1}{c_1F_m(x,t)}-\frac{2\Delta\eta(x)}{c_1} -\frac{4}{c_1F_m(x,t)}
\nabla\eta(x)\cdot\nabla F_m(x,t) \\
&\le &  \frac{\eta C_1}{c_1C_0}-\frac{2\Delta\eta(x)}{c_1}
-\frac{4}{c_1F_m(x,t)} \nabla\eta(x)\cdot\nabla F_m(x,t)\end{eqnarray*}
Since
$(x,t)$ is a maximum for $\eta F_m$ we have
$$0=\nabla(\eta(x) F_m(x,t))=\nabla\eta(x) F_m(x,t)+\eta(x)\nabla F_m(x,t),$$
so that
$$\frac{\nabla\eta(x)}{\eta(x)}=-\frac{\nabla F_m(x,t)}{F_m(x,t)}.$$
Plugging this in gives
$$\eta(x) F_m(x,t)\le
\frac{C_1}{c_1C_0}-\frac{2\Delta\eta(x)}{c_1}+4\frac{|\nabla\eta(x)|^2}{c_1\eta(x)}+\eta
C_0.$$

Of course, the gradient and Laplacian of $\eta$ are taken at the
point $(x,t)$.  Thus, because of the properties of $\eta$ given in
Proposition~\ref{prop:parabolic-cutoff-function}, it immediately follows that $\eta F_m(x,t)$ is
bounded by a constant depending
 only on $K,n,\alpha,r,c_1,C_0,C_1$, and as we have already seen, $c_1,C_0,C_1$ depend only on
 $K,\alpha,A_1,\ldots,A_m$.

Now suppose that $s\left\{  \max\left\{  m-l+1,0\right\}
\right\}=1.$ Again we compute the evolution inequality for $\eta
F_{m}$. The result is $$\left(  \frac{\partial}{\partial
t}-\Delta\right) \left(  \eta F_{m}\right) \leq\eta\left(
-\frac{c_1}{t} (F_{m}-C_0)^{2}+\frac{C_1}{t}\right) -\Delta\eta\cdot
F_m-2\nabla\eta\cdot\nabla F_m.$$ Thus,  using the maximum principle
as before, we have
$$\left(  \frac{\partial}{\partial t}-\Delta\right)\eta  F_{m}(x,t)\ge 0.$$
Hence,
$$\frac{\eta(x) c_1(F_m(x,t)-C_0)^2}{t}\le
\frac{\eta(x)C_1}{t}-\Delta\eta(x)F_m(x,t)-2\nabla\eta(x)\cdot\nabla F_m(x,t).$$
Using the assumption that $F_m(x,t)\ge 2C_0$ as before, and rewriting the last
term as before, we have
$$\eta F_m(x,t)\le
\frac{\eta(x)C_1}{c_1C_0}-\frac{2t\Delta\eta(x)}
{c_1}+\frac{4t|\nabla\eta(x)|^2}{c_1\eta(x)}+\eta C_0.$$ The right-hand side is
bounded by a constant depending only on $K,n,\alpha,r,c_1,C_0,C_1$. We conclude
that in all cases $\eta F_{m}$ is bounded by a constant depending only on
$K,n,\alpha,r,c_1,C_0,C_1$, and hence on $K,n,\alpha,r,A_1,\ldots,A_m$.


This proves that for any $y\in B(p,0,r/2)$, the value
 $\eta F_m(x,t)$ is bounded by a
constant $A_{m+1}$ depending only on $(m+1,l,K,n,\alpha,r)$ for all
$(x,t)\in B(y,0,r/2)\times [0,T]$. Since $\eta(y)=1$, for all $0\le t\le T$ we have
$$t^{\max\{m+1-l,0\}}|\nabla^{m+1}\Rm(y,t)|^2\le F_m(y,t)=\eta(y)F_m(y,t)\le A_{m+1}.$$



This completes the inductive proof that the result holds for
$k=m+1$ and hence establishes Theorem~\ref{thm:shi-with-initial-derivatives}, modulo the proof of Claim~\ref{claim:function-evolution-bound}.
\end{proof}

\begin{remark}[Shi derivative variant comparison]\label{rem:shi-derivative-variant}
\uses{thm:shi-with-initial-derivatives, thm:shi-first-derivative}

The case $l=0$ of Theorem~\ref{thm:shi-with-initial-derivatives} is Shi's theorem
(Theorem~\ref{thm:shi-first-derivative}).
\end{remark}

\begin{corollary}[C-infinity curvature bounds]\label{cor:ricci-flow-curvature-bounds}
\uses{thm:shi-with-initial-derivatives}
\label{Cinftybd}

Suppose that $(M,g(t)),\ 0\le t\le T$, is a Ricci flow
with $(M,g(t))$ complete and $T<\infty$, with $\Rm(x,0)$ bounded in
the $C^\infty$-topology independently of $x$, and with $|\Rm(x,t)|$ bounded
independently of $(x,t)\in M\times[0,T]$. Then $\Rm(x,t)$ is bounded in the
$C^\infty$-topology independently of $(x,t)\in M\times[0,T]$.
\end{corollary}




Now we turn to the proof of Claim~\ref{claim:function-evolution-bound}.

\begin{proof}
  \uses{prop:parabolic-cutoff-function, def:ricci-flow}
\sketch In this argument fix
$(x,t)\in B(p,0,3r/4)\times (0,T]$ and we drop $(x,t)$ from the
notation. Recall that by Equations (7.4a) and (7.4b) on p. 229 of
\cite{ChowKnopf}  we have
\begin{align}
\frac{\partial}{\partial t}\left\vert \nabla^{\ell
}\Rm\right\vert ^{2} & \leq\Delta\left\vert
\nabla^{\ell }\Rm\right\vert ^{2} -2\left\vert
\nabla^{\ell +1}\Rm\right\vert ^{2}+\sum_{i=0} ^{\ell
}c_{\ell,j}\left\vert \nabla^{i}\Rm\right\vert
\left\vert \nabla^{\ell -i}\Rm\right\vert \left\vert
\nabla^{\ell } \Rm\right\vert, \label{curderfor}
\end{align}
where the constants $c_{\ell,j}$ depend only on $\ell$ and $j$.


Hence, setting $m_l= \max\left\{ m+1-l,0\right\}$ and denoting
$c_{m+1,i}$ by $\tilde c_i$, we have

\begin{eqnarray}
\lefteqn{\frac{\partial}{\partial t}\left(  t^{m_l}\left\vert
\nabla^{m+1}\Rm\right\vert ^{2}\right)
 \leq\Delta\left(  t^{m_l }\left\vert \nabla
^{m+1}\Rm\right\vert ^{2}\right)  -2t^{m_l }
\left\vert \nabla^{m+2}\Rm\right\vert ^{2}}\label{1steqn}\\
& & +t^{m_l  }\sum_{i=0}^{m+1}\tilde c_{i}\left\vert \nabla
^{i}\Rm\right\vert \left\vert
\nabla^{m+1-i}\Rm \right\vert \left\vert
\nabla^{m+1}\Rm\right\vert
 +m_l  t^{m_l -1}\left\vert \nabla^{m+1}\Rm\right\vert ^{2}\nonumber\\
&  \leq & \Delta\left(  t^{m_l }\left\vert \nabla
^{m+1}\Rm\right\vert ^{2}\right)  -2t^{m_l }
\left\vert \nabla^{m+2}\Rm\right\vert ^{2} +(\tilde
c_{0}+\tilde
c_{m+1})t^{m_l}\left\vert\Rm\right\vert\left\vert\nabla^{m+1}\Rm
\right\vert^2\nonumber\\ & &+ t^{m_l }\sum_{i=1}^{m}\tilde
c_{i}\left\vert \nabla
^{i}\Rm\right\vert \left\vert \nabla^{m+1-i}\Rm%
\right\vert \left\vert \nabla^{m+1}\Rm\right\vert
+m_l  t^{m_l -1}\left\vert \nabla^{m+1}\Rm\right\vert
^{2}.\nonumber
\end{eqnarray}

Using the inductive hypothesis, Inequality~(\ref{strongineq}), there is a constant
$A<\infty$ depending only on $K,\alpha,A_1,\ldots,A_m$ such that
$$
\sum_{i=1}^{m}\tilde c_{i}\left\vert \nabla
^{i}\Rm\right\vert \left\vert
\nabla^{m+1-i}\Rm \right\vert \leq At^{-m_l/2}.
$$
Also, let $c=\tilde c_{0}+\tilde c_{m+1}$ and define a new  constant
$B$
by%\note{I think this should be independent of $K$.}
$$B=c(\alpha+K)+m_l.$$ Then, since $t\le T\le \alpha/K$ and $m_l\ge 0$, we have
$$((\tilde c_{0}+\tilde c_{m+1})t\left\vert\Rm\right\vert +m_l)t^{m_l-1}\le
\frac{Bt^{m_l}}{t^{s(m_l)}}.$$ Putting this together allows us to
rewrite Inequality~(\ref{1steqn}) as
\begin{align*}
\frac{\partial}{\partial t}\left(  t^{m_l}\left\vert
\nabla^{m+1}\Rm\right\vert ^{2}\right)
 \leq &\Delta\left(  t^{m_l }\left\vert \nabla
^{m+1}\Rm\right\vert ^{2}\right)  -2t^{m_l }
\left\vert \nabla^{m+2}\Rm\right\vert ^{2}\\
 &
+At^{m_l /2}\left\vert \nabla^{m+1} \Rm\right\vert
+\left(ct\left\vert \Rm\right\vert +m_l \right)
t^{m_l -1}\left\vert \nabla^{m+1}\Rm\right\vert
^{2}\\
 \leq &\Delta\left(  t^{m_l }\left\vert \nabla
^{m+1}\Rm\right\vert ^{2}\right)  -2t^{m_l }
\left\vert \nabla^{m+2}\Rm\right\vert ^{2}\\
& +
\frac{B}{t^{s(m_l)}}t^{m_l}\left\vert\nabla^{m+1}\Rm\right\vert^2
+ At^{m_l/2}\left\vert\nabla^{m+1}\Rm\right\vert.
\end{align*}
Completing the square gives
\begin{align*}
\frac{\partial}{\partial t}\left(  t^{m_l}\left\vert
\nabla^{m+1}\Rm\right\vert ^{2}\right)
 &
\leq\Delta\left( t^{m_l }\left\vert \nabla
^{m+1}\Rm\right\vert ^{2}\right) -2t^{m_l  }
\left\vert \nabla^{m+2}\Rm\right\vert ^{2}\\
&  +(B+1)t^{m_l  -s(m_l)}\left\vert \nabla^{m+1}
\Rm\right\vert ^{2}+\frac{A^2}{4}t^{s(m_l)}.
\end{align*}

Let $\hat{m}_l = \max\left\{  m-l,0\right\} $. From
(\ref{curderfor}) and the induction hypothesis, there is a constant
$D$, depending on $K,\alpha,A_1,\ldots,A_m$ such that

\begin{eqnarray*}
\frac{\partial}{\partial t}\left(  t^{\hat{m}_l }\left\vert
\nabla^{m}\Rm\right\vert ^{2}\right) & \leq &
\Delta\left( t^{\hat{m}_l} \left\vert
\nabla^{m}\Rm\right\vert ^{2}\right) -2t^{\hat{m}_l
}\left\vert
\nabla^{m+1}\Rm\right\vert ^{2}\\
& & +\hat{m}_l  t^{\hat{m}_l  -1}\left\vert
\nabla^{m}\Rm\right\vert ^{2}+D.
\end{eqnarray*}

Now, defining new constants $\widetilde B=B+1$ and $\widetilde
A=A^2/4$ we have

\begin{eqnarray*}
\lefteqn{ \left(  \frac{\partial}{\partial t}-\Delta\right)F_m=
\left( \frac{\partial}{\partial t}-\Delta\right)  \left[ \left(
C+t^{\hat{m}_l  }\left\vert \nabla^{m}\Rm \right\vert
^{2}\right)  t^{m_l }\left\vert
\nabla^{m+1}\Rm\right\vert ^{2}\right]  \leq}   \\
& & \left(  C+t^{\hat{m}_l  }\left\vert \nabla^{m}
\Rm\right\vert ^{2}\right)  \left( -2t^{m_l
}\left\vert \nabla^{m+2} \Rm \right\vert ^{2}
+\frac{\widetilde B}{t^{s\{m_l\}}} t^{m_l}\left\vert \nabla^{m+1}
\Rm \right\vert ^{2}+\widetilde At^{s(m_l)} \right) \\
& &  +\left( -2t^{\hat{m}_l }\left\vert
\nabla^{m+1}\Rm \right\vert ^{2} +\hat{m}_l
t^{\hat{m}_l  -1}\left\vert \nabla^{m}\Rm\right\vert
^{2}+D \right) t^{m_
l} \left\vert \nabla^{m+1} \Rm\right\vert ^{2}\\
& &  -2t^{\hat{m}_l  +m_l } \nabla\left(\left\vert
\nabla^{m}\Rm\right\vert ^{2}\right)\cdot
\nabla\left(\left\vert \nabla^{m+1}\Rm\right\vert
^{2}\right). \end{eqnarray*}

Since $C\ge 4t^{\hat
m_l}\left\vert\nabla^m\Rm\right\vert^2$, this implies

\begin{eqnarray}
\lefteqn{ \left(  \frac{\partial}{\partial t}-\Delta\right) F_m \leq
-10t^{\hat{m}_l  +m_l} \left\vert
\nabla^{m}\Rm\right\vert ^{2}\left\vert
\nabla^{m+2}\Rm\right\vert ^{2}}\label{2ndeqn}\\
 & & -8t^{\hat{m}_l  +m_l  }\left\vert
\nabla^{m}\Rm\right\vert \left\vert \nabla^{m+1}%
\Rm\right\vert ^{2}\left\vert \nabla^{m+2}\Rm%
\right\vert  -2t^{\hat m_l  +m_l  }\left\vert
\nabla^{m+1}\Rm\right\vert ^{4}\nonumber\\
 & & +\left(  C+t^{\hat m_l}\left\vert \nabla^{m}
\Rm\right\vert ^{2}\right)  \left( \widetilde
Bt^{m_l-s(m_l) }\left\vert \nabla^{m+1}\Rm\right\vert
^{2}+\widetilde At^{s(m_l)}\right) \nonumber\\
 & & +\left(\hat m_l  t^{\hat{m}_l -1}\left\vert
\nabla^{m}\Rm\right\vert ^{2}+D\right) t^{m_l
}\left\vert \nabla^{m+1}\Rm \right\vert
^{2}.\nonumber
\end{eqnarray}

Now we can write the first three terms on the right-hand side of
Inequality~(\ref{2ndeqn}) as \begin{equation}\label{3rdeqn} -t^{\hat
m_l+m_l}\left(\sqrt{10}\left\vert\nabla^{m+2}\Rm\right\vert
\left\vert\nabla^m\Rm\right\vert+\frac{4}{\sqrt{10}}\left\vert\nabla^{m+1}
\Rm\right\vert^2\right)^2-\frac{2}{5}t^{\hat
m_l+m_l}\left\vert\nabla^{m+1}\Rm\right\vert^4.\end{equation}
In addition we have
\begin{equation}\label{4theqn}
C+t^{\hat m_l}\left\vert \nabla^{m} \Rm\right\vert
^{2}  \leq  C+A_m^2.
\end{equation}

Let us set $\widetilde D=\max(\alpha/K,1)D$. If $\hat m_l=0$, then
\begin{equation}\label{5theqn}
\hat m_lt^{\hat
m_l-1}\left\vert\nabla^m\Rm\right\vert^2+D =D\le
\frac{\widetilde D}{t^{s(m_l)}}=\hat m_lA_m^2+D\le \frac{\hat
m_lA_m^2+\widetilde D}{t^{s(m_l)}}.\end{equation}  On the other
hand, if $\hat m_l>0$, then $s(\hat m_l)=s(m_l)=1$ and hence
$$\hat m_lt^{\hat m_l-1}\left\vert\nabla^m\Rm\right\vert^2+D\le
\frac{1}{t^{s(m_l)}}\hat m_lA_m^2+D\le \frac{\hat
m_lA_m^2+\widetilde D}{t^{s(m_l)}}.$$


Since $\hat m_l =m_l-s(m_l)$,
Inequalities~(\ref{3rdeqn}),~(\ref{4theqn}), and~(\ref{5theqn}) then
allow us  rewrite Inequality~(\ref{2ndeqn}) as
\begin{eqnarray*}
\lefteqn{ \left(  \frac{\partial}{\partial t}-\Delta\right)F_m
 \leq   -\frac{2}{5t^{s( m_l)  }
 }t^{2m_l}\left\vert\nabla^{m+1}\Rm\right\vert^4}\\
& &  +(C+A_m^2)\left(\frac{\widetilde B}{t^{s(m_l)}}
t^{m_l}\left\vert\nabla^{m+1}\Rm\right\vert^2+\widetilde
At^{s(m_l)}\right) +\frac{\hat m_lA_m^2+\widetilde
D}{t^{s(m_l)}}t^{m_l}\left\vert\nabla^{m+1}\Rm
\right\vert^2.
\end{eqnarray*}

Setting
$$B'=(C+A_m^2)\widetilde B+(\hat m_lA_m^2+\widetilde D),$$
and $A'=\widetilde A(C+A_m^2)$ we have
\begin{eqnarray*}
\left(  \frac{\partial}{\partial t}-\Delta\right) F_m  & \leq &
-\frac{2}{5t^{s( m_l)  }
 }\left(t^{m_l}\left\vert\nabla^{m+1}\Rm\right\vert^2\right)^2 \\
 & & +
\frac{B'}{t^{s(m_l)}}t^{m_l}\left\vert\nabla^{m+1}\Rm\right\vert^2
+A't^{s(m_l)}.
\end{eqnarray*}
We rewrite this as
\begin{eqnarray*}
\left(  \frac{\partial}{\partial t}-\Delta\right) F_m & \leq &
-\frac{2}{5t^{s(m_l)}}\left(t^{m_l}\left\vert\nabla^{m+1}\Rm
\right\vert^2   -\frac{5B'}{4}\right)^2
\\ & & +\frac{5(B')^2}{8t^{s(m_l)}} +A't^{s(m_l)},
\end{eqnarray*}
and hence
\begin{equation*}
\left(  \frac{\partial}{\partial t}-\Delta\right) F_m \leq
-\frac{2}{5t^{s(m_l)}}\left(t^{m_l}\left\vert\nabla^{m+1}\Rm
\right\vert^2-B''\right)^2  +\frac{A''}{t^{s(m_l)}}
\end{equation*}
where the constants $B''$ and $A''$ are defined by $ B''=5 B'/4$ and
$$A''=(\max\{\alpha/K,1\})^2+5( B')^2/8.$$ (Recall that $t\le T\le
\alpha/K$.) Let
$$Y=(C+t^{\hat
m_l}\left\vert\nabla^m\Rm\right\vert^2).$$ (Notice
that $Y$ is not a constant.) Of course, by definition
$$F_m=Yt^{m_l}|\nabla^{m+1}\Rm|^2.$$ Then the previous inequality
becomes
\begin{equation*}
\left(\frac{\partial}{\partial t}-\Delta\right)  F_m \leq
-\frac{2}{5t^{s(m_l)}Y^2}\left(Yt^{m_l}\left\vert\nabla^{m+1}\Rm
\right\vert^2-B'' Y\right)^2
 +\frac{A''}{t^{s(m_l)}}
\end{equation*}
Since $C\le Y\le 5C/4$ we have
\begin{eqnarray*}
\left(  \frac{\partial}{\partial t}-\Delta\right)  F_m \leq
-\frac{32}{125t^{s(m_l)}C^2}\left(F_m-B''
Y\right)^2+\frac{A''}{t^{s(m_l)}}
\end{eqnarray*}
At any point where $F_m\ge 5CB''/4$,  the last inequality gives
$$\left(  \frac{\partial}{\partial t}-\Delta\right)  F_m \leq
-\frac{32}{125t^{s(m_l)}C^2}\left(F_m-5CB''/4\right)^2+\frac{A''}{t^{s(m_l)}}.$$
At any point where $F_m\le 5CB''/4$,  since $F_m\ge 0$ and $0\le
B''Y\le 5CB''/4$, we have $(F_m-B''Y)^2\le 25C^2(B'')^2/16$, so that
$$-\frac{32}{125t^{s(m_l)}C^2}\left(F_m-5CB''/4\right)^2\ge
-2(B'')^2/5t^{s(m_l)}.$$ Thus, in this case we have
$$\left(\frac{\partial}{\partial t}-\Delta\right)  F_m \leq
\frac{A''}{t^{s(m_l)}}\le
-\frac{32}{125t^{s(m_l)}C^2}\left(F_m-5CB''/4\right)^2+\frac{A''+2(B'')^2/5}{t^{s(m_l)}}.$$


These two cases together prove Claim~\ref{claim:function-evolution-bound}.
\end{proof}




\section{Generalized Ricci flows}

In this section we introduce a generalization of the Ricci flow
equation. The generalization does not involve changing the PDE that
gives the flow. Rather it allows for the global topology of
space-time to be different from a product.
\subsection{Space-time}

There are two basic ways to view an $n$-dimensional Ricci flow: (i)
as a one-parameter family of metrics $g(t)$ on a fixed smooth
$n$-dimensional manifold $M$, and (ii) as a partial metric (in the
horizontal directions) on the $(n+1)$-dimensional manifold $M\times
I$. We call the latter $(n+1)$-dimensional manifold {\sl
space-time} and the
horizontal slices are the {\em time-slices}. In defining the
generalized Ricci flow, it is the second approach that we
generalize.

\begin{definition}[Space-time]\label{def:generalized-spacetime}
\lean{MorganTianLib.GeneralizedSpaceTimeLocalChart, MorganTianLib.GeneralizedSpaceTime}
\leanok
By {\em space-time} we mean a smooth $(n+1)$-dimensional manifold
${\mathcal M}$ (possibly with boundary), equipped with a smooth
function ${\bf t}\colon {\mathcal M}\to \Ar$, called {\em time} and
a smooth vector field $\chi$ subject to the following axioms:
\begin{enumerate}
\item The image of ${\bf t}$ is an interval $I$ (possibly infinite) and
the boundary of ${\mathcal M}$ is the preimage under ${\bf t}$ of
$\partial I$.
\item For each $x\in {\mathcal M}$ there is an open neighborhood $U\subset
{\mathcal M}$ of $x$ and a diffeomorphism $f\colon  V\times J\to U$,
where $V$ is an open subset in $\Ar ^n$ and $J$ is an interval with
the property that (i) ${\bf t}$ is the composition of $f^{-1}$
followed by the projection onto the interval $J$ and (ii) $\chi$ is
the image under $f$ of the unit vector field in the positive
direction tangent to the foliation by the lines $\{v\}\times J$ of
$V\times J$.
\end{enumerate}
\end{definition}

 Notice that it follows
that $\chi({\bf t})=1$.




\begin{definition}[Time-slices and horizontal metric]\label{def:time-slices-horizontal}
\uses{def:generalized-spacetime}
\lean{MorganTianLib.GeneralizedSpaceTime.timeSlice, MorganTianLib.GeneralizedSpaceTime.HorizontalTangentSpace, MorganTianLib.GeneralizedSpaceTime.horizontalProjectionAt, MorganTianLib.GeneralizedSpaceTime.horizontal_finrank_add_one, MorganTianLib.GeneralizedSpaceTimeLocalChart.preimage_timeSlice, MorganTianLib.GeneralizedSpaceTime.boundary_connectedComponent_time_eq, MorganTianLib.GeneralizedSpaceTime.HorizontalMetric, MorganTianLib.GeneralizedSpaceTime.HorizontalMetric.horizontalInner, MorganTianLib.GeneralizedSpaceTime.timeDifferential_timeVector}
\leanok

The {\em time-slices}
of space-time are the level sets ${\bf t}$. These form a
codimension-one foliation of ${\mathcal M}$. For each $t\in I$ we
denote by $M_t\subset {\mathcal M}$ the $t$ time-slice, that is to
say ${\bf t}^{-1}(t)$. Notice that each boundary component of
${\mathcal M}$ is contained in a single time-slice. The {\em
horizontal distribution}, ${\mathcal H}T{\mathcal M}$ is the
distribution tangent to this foliation. A {\em horizontal
metric} on
space-time is a smoothly varying positive definite inner product on
${\mathcal H}T{\mathcal M}$.
\end{definition}

Notice that a horizontal metric on space-time induces an ordinary
Riemannian metric on each time-slice. Conversely, given a Riemannian
metric on each time-slice $M_t$, the condition that they fit
together to form a horizontal metric on space-time is that they vary
smoothly on space-time. We define the curvature of a horizontal
metric $G$ to be the section of the dual of the symmetric square of
$\wedge^2{\mathcal H}T{\mathcal M}$ whose value at each point $x$ with ${\bf
t}(x)=t$ is the usual Riemannian curvature tensor of the induced
metric on $M_t$  at the point $x$. This is a smooth section of $\Sym^2(\wedge^2{\mathcal H}T^*{\mathcal M})$.
 The Ricci curvature and the scalar curvature of a horizontal metric
 are given in the usual way from its Riemannian curvature.
 The Ricci curvature is a smooth section of $\Sym^2({\mathcal
 H}T^*{\mathcal M})$ while the scalar curvature is a smooth function on ${\mathcal M}$.

\subsection{The generalized Ricci flow equation}



Because of the second condition in the definition of space-time, the
vector field $\chi$ preserves the horizontal foliation and hence the
horizontal distribution. Thus, we can form the Lie derivative of a
horizontal metric with respect to $\chi$.

\begin{definition}[Generalized Ricci flow]\label{def:generalized-ricci-flow}
\uses{def:generalized-spacetime, def:time-slices-horizontal}
\lean{MorganTianLib.GeneralizedRicciFlow, MorganTianLib.productSpaceTimeTime, MorganTianLib.productSpaceTimeTimeVector, MorganTianLib.productSpaceTimeTime_mfderiv_timeVector_formula, MorganTianLib.productSpaceTimeSlice, MorganTianLib.ProductGeneralizedRicciFlow, MorganTianLib.ProductGeneralizedRicciFlow.of_isRicciFlowOn, MorganTianLib.ProductGeneralizedRicciFlow.to_isRicciFlowOn, MorganTianLib.ProductGeneralizedRicciFlow.metricSection, MorganTianLib.ProductGeneralizedRicciFlow.metricSection_smooth, MorganTianLib.ProductGeneralizedRicciFlow.metricSection_ricci_equation}
\leanok

An {\em $n$-dimensional generalized Ricci flow} consists of a space-time ${\mathcal M}$ that is
$(n+1)$-dimensional and a horizontal metric $G$ satisfying the
generalized Ricci flow equation:
$${\mathcal L}_\chi(G)=-2\Ric(G).$$
\end{definition}


\begin{remark}[Generalized Ricci flow local behavior]\label{rem:generalized-ricci-flow-local}
\uses{def:generalized-ricci-flow}
\lean{MorganTianLib.GeneralizedRicciFlow.LocalOrdinaryRicciFlowAt, MorganTianLib.GeneralizedRicciFlow.localOrdinaryRicciFlowAt}

Let $({\mathcal M},G)$ be a generalized Ricci flow and let $x\in
{\mathcal M}$. Pulling $G$ back to  the local coordinates $V\times
J$ defined near any point gives a one-parameter family of metrics
$(V,g(t)),\  t\in J$, satisfying the usual Ricci flow equation. It
follows that all the usual evolution formulas for Riemannian
curvature, Ricci curvature, and scalar curvature hold in this more
general context.
\end{remark}

Of course, any ordinary Ricci flow is a generalized Ricci flow where
space-time is a product $M\times I$ with time being the projection
to $I$ and $\chi$ being the unit vector field in the positive
$I$-direction.


\subsection{More definitions for generalized Ricci flows}

\begin{definition}[Compatible embedding in space-time]\label{def:compatible-embedding}
\uses{def:generalized-spacetime}
\lean{MorganTianLib.GeneralizedSpaceTime.CompatibleEmbedding, MorganTianLib.GeneralizedSpaceTime.CompatibleEmbedding.image, MorganTianLib.GeneralizedSpaceTime.CompatibleEmbedding.time_toFun, MorganTianLib.GeneralizedSpaceTime.CompatibleEmbedding.toFun_mem_timeSlice, MorganTianLib.GeneralizedSpaceTime.CompatibleEmbedding.image_subset_time_preimage, MorganTianLib.GeneralizedSpaceTime.CompatibleEmbedding.time_image, MorganTianLib.GeneralizedSpaceTime.CompatibleEmbedding.isInteriorPoint_toFun_Ioo, MorganTianLib.GeneralizedSpaceTime.CompatibleTimeSliceEmbedding, MorganTianLib.GeneralizedSpaceTime.CompatibleTimeSliceEmbedding.toFun_center, MorganTianLib.GeneralizedSpaceTime.CompatibleTimeSliceEmbedding.image_inter_timeSlice, MorganTianLib.GeneralizedSpaceTime.CompatibleTimeSliceEmbedding.eqOn_Ioo, MorganTianLib.GeneralizedSpaceTime.CompatibleTimeSliceEmbedding.eqOn_Icc_of_closure, MorganTianLib.GeneralizedSpaceTime.CompatibleTimeSliceEmbedding.eqOn_Icc_of_isClosed, MorganTianLib.GeneralizedSpaceTime.CompatibleTimeSliceEmbedding.eqOn_Icc_of_embedding, MorganTianLib.GeneralizedSpaceTime.CompatibleEmbedding.affineTimeReparam_isEmbedding, MorganTianLib.GeneralizedSpaceTime.FlowLine, MorganTianLib.GeneralizedSpaceTime.CompatibleTimeSliceDiffeomorphism, MorganTianLib.GeneralizedSpaceTime.CompatibleTimeSliceDiffeomorphism.image_eq}

Let ${\mathcal M}$ be a space-time. Given a space $C$ and an
interval $I\subset \Ar$  we say that an embedding $C\times I\to
{\mathcal M}$ is {\em compatible with} the time and the vector field
if: (i) the restriction of ${\bf t}$ to the image agrees with the
projection onto the second factor and (ii) for each $c\in C$ the
image of $\{c\}\times I$ is the integral curve for the vector field
$\chi$.
 If in addition $C$ is a subset of $M_{t}$
we require that $t\in I$ and that the map $C\times\{t\}\to M_{t}$ be
the identity. Clearly, by the uniqueness of integral curves for vector fields,
 two such embeddings agree on their common interval of definition, so that,
given $C\subset M_{t}$ there is a maximal interval $I_C$ containing
$t$ such that such an embedding, compatible with time and the vector field,
 is defined on $C\times I$. In the
special case when $C=\{x\}$ for a point $x\in M_{t}$ we say that
such an embedding is {\em the flow line} through $x$. The embedding
of the maximal interval through $x$ compatible with time and the vector
field $\chi$ is called {\em the domain of definition} of the flow
line through $x$. For a more general subset $C\subset M_{t}$ there
is an embedding $C\times I$ compatible with time and the vector
field $\chi$ if an only if for every $x\in C$, $I$ is contained in
the domain of definition of the flow line through $x$.
\end{definition}


\begin{definition}[Regular and singular times]\label{def:regular-singular-time}
\uses{def:compatible-embedding, def:time-slices-horizontal}
\lean{MorganTianLib.GeneralizedSpaceTime.IsRegularTime, MorganTianLib.GeneralizedSpaceTime.IsSingularTime, MorganTianLib.GeneralizedSpaceTime.isSingularTime_iff, MorganTianLib.GeneralizedSpaceTime.IsRegularTime.mem_timeRange, MorganTianLib.GeneralizedSpaceTime.IsSingularTime.mem_timeRange, MorganTianLib.GeneralizedSpaceTime.IsRegularTime.not_isSingularTime, MorganTianLib.GeneralizedSpaceTime.regular_or_singular, MorganTianLib.GeneralizedSpaceTime.IsInitialTime, MorganTianLib.GeneralizedSpaceTime.HasInitialTimeNegInfinity, MorganTianLib.GeneralizedSpaceTime.isInitialTime_unique, MorganTianLib.GeneralizedSpaceTime.IsInitialTime.le_time, MorganTianLib.GeneralizedSpaceTime.IsInitialTime.exists_time_mem_Ico, MorganTianLib.GeneralizedSpaceTime.exists_isInitialTime_of_bddBelow, MorganTianLib.GeneralizedSpaceTime.hasInitialTimeNegInfinity_of_not_bddBelow, MorganTianLib.GeneralizedSpaceTime.exists_isInitialTime_or_hasInitialTimeNegInfinity, MorganTianLib.GeneralizedSpaceTime.timeRange_isOpen_of_all_regular, MorganTianLib.GeneralizedSpaceTime.IsRegularTime.not_isInitialTime, MorganTianLib.GeneralizedSpaceTime.IsInitialTime.not_hasInitialTimeNegInfinity}

We say that $t$ is a {\em regular time} if there is $\epsilon>0$ and
a diffeomorphism $M_{t}\times (t-\epsilon,t+\epsilon)\to {\bf
t}^{-1}((t-\epsilon,t+\epsilon))$ compatible with time
and the vector field. A time is {\em singular} if it is not regular.
Notice that if all times are regular, then space-time is a product
$M_{t}\times I$ with ${\bf t}$ and $\chi$ coming from the second
factor. If the image ${\bf t}({\mathcal M})$ is an interval $I$
bounded below, then the {\em initial time} for the flow is the
greatest lower bound for $I$. If $I$ includes $(-\infty,A]$ for some
$A$, then the initial time for the generalized Ricci flow is
$-\infty$.
\end{definition}



\begin{definition}[Scaling and translation of generalized Ricci flows]\label{def:scaling-translation}
\uses{def:generalized-ricci-flow}
\lean{MorganTianLib.parabolicTimeOrderIso, MorganTianLib.parabolicTimeOrderIso_image_Icc, MorganTianLib.parabolicTimeOrderIso_image_Ioo, MorganTianLib.parabolicTimeOrderIso_image_Ico, MorganTianLib.parabolicTimeOrderIso_image_Ioc, MorganTianLib.parabolicTimeOrderIso_image_Ici, MorganTianLib.parabolicTimeOrderIso_image_Iic, MorganTianLib.parabolicTimeOrderIso_comp, MorganTianLib.parabolicTimeOrderIso_image_image, MorganTianLib.parabolicTimeOrderIso_preimage_image, MorganTianLib.parabolicTimeOrderIso_image_preimage, MorganTianLib.parabolicTimeOrderIso_product_image_image, MorganTianLib.GeneralizedSpaceTime.CompatibleEmbedding.affineTimeReparam, MorganTianLib.GeneralizedSpaceTime.CompatibleEmbedding.affineTimeReparam_image, MorganTianLib.GeneralizedSpaceTime.CompatibleEmbedding.affineTimeReparam_time, MorganTianLib.GeneralizedSpaceTime.CompatibleEmbedding.affineTimeReparam_time_image, MorganTianLib.GeneralizedSpaceTime.CompatibleEmbedding.isIntegralCurveOn_affineTimeReparam, MorganTianLib.GeneralizedSpaceTime.CompatibleEmbedding.affineTimeReparam_isEmbedding, MorganTianLib.GeneralizedSpaceTime.scaledTime, MorganTianLib.GeneralizedSpaceTime.translatedTime, MorganTianLib.GeneralizedSpaceTime.range_scaledTime, MorganTianLib.GeneralizedSpaceTime.range_translatedTime, MorganTianLib.GeneralizedSpaceTime.frontier_range_scaledTime, MorganTianLib.GeneralizedSpaceTime.frontier_range_translatedTime, MorganTianLib.GeneralizedSpaceTime.scaledTime_preimage_frontier_range, MorganTianLib.GeneralizedSpaceTime.translatedTime_preimage_frontier_range, MorganTianLib.GeneralizedSpaceTime.scaledTimeVector, MorganTianLib.GeneralizedSpaceTime.affineTimeChange, MorganTianLib.GeneralizedSpaceTime.affineTimeChange_time, MorganTianLib.GeneralizedSpaceTime.affineTimeChange_timeVector, MorganTianLib.GeneralizedSpaceTime.affineTimeChange_range, MorganTianLib.GeneralizedSpaceTime.affineTimeChange_frontier_range, MorganTianLib.GeneralizedSpaceTime.affineTimeChange_preimage_frontier_range, MorganTianLib.GeneralizedSpaceTime.affineTimeChange_isHorizontal_iff, MorganTianLib.GeneralizedSpaceTime.HorizontalMetric.affineTimeChange, MorganTianLib.GeneralizedSpaceTime.HorizontalMetric.affineTimeChange_inner, MorganTianLib.GeneralizedSpaceTime.scaledTimeDifferential, MorganTianLib.GeneralizedSpaceTime.translatedTimeDifferential, MorganTianLib.GeneralizedSpaceTime.scaledTimeDifferential_eq, MorganTianLib.GeneralizedSpaceTime.translatedTimeDifferential_eq, MorganTianLib.GeneralizedSpaceTime.HorizontalMetric.inner_self_nonneg, MorganTianLib.GeneralizedSpaceTime.HorizontalMetric.constScale, MorganTianLib.GeneralizedSpaceTime.HorizontalMetric.constScale_inner, MorganTianLib.GeneralizedSpaceTimeLocalChart.realizesHorizontalMetric_constScale, MorganTianLib.GeneralizedSpaceTimeLocalChart.affineTimeChange_spatialPlaque_eq, MorganTianLib.GeneralizedSpaceTimeLocalChart.affineTimeChange_spatialDifferential, MorganTianLib.parabolicRescaledMetricFamily, MorganTianLib.parabolicRescaledMetricFamily_metricInner, MorganTianLib.isRicciFlowEquationOn_parabolicRescaledMetricFamily, MorganTianLib.GeneralizedSpaceTimeLocalChart.isRicciFlowEquationOn_parabolicRescaledMetricFamily, MorganTianLib.GeneralizedRicciFlow.affineTimeChange, MorganTianLib.GeneralizedSpaceTime.CompatibleEmbedding.affineTimeReparam_parabolicTimeOrderIso, MorganTianLib.GeneralizedSpaceTime.HorizontalMetric.horizontalCurveSpeed_constScale, MorganTianLib.GeneralizedSpaceTime.HorizontalMetric.horizontalCurveLength_constScale, MorganTianLib.GeneralizedSpaceTime.HorizontalMetric.timeSliceEDist_constScale, MorganTianLib.GeneralizedSpaceTime.HorizontalMetric.timeSliceEDist_constScale_lt_ofReal_iff, MorganTianLib.isRicciFlowOn_parabolicRescaledMetricFamily_of_smooth, MorganTianLib.parabolicRescaledMetricFamily_sourceTime_mem}

Suppose that $({\mathcal M},G)$ is a generalized Ricci flow and that
$Q>0$ is a positive constant. Then we can define a new generalized
Ricci flow by setting $G'=QG$, ${\bf t'}=Q{\bf t}$ and
$\chi'=Q^{-1}\chi$. It is easy to see that the result still
satisfies the generalized Ricci flow equation. We denote this new
generalized Ricci flow by $(Q{\mathcal M},QG)$ where the changes in
${\bf t}$ and $\chi$ are denoted by the factor of $Q$ in front of
${\mathcal M}$.

It is also possible to translate a generalized solution $({\mathcal M},G)$ by
replacing the time function ${\bf t}$ by ${\bf t'}={\bf t}+a$ for any constant
$a$, leaving $G$ and $\chi$ unchanged.
\end{definition}

\begin{definition}[Parabolic neighborhoods]\label{def:parabolic-neighborhoods}
\uses{def:generalized-spacetime, def:generalized-ricci-flow, def:compatible-embedding}
\lean{MorganTianLib.GeneralizedSpaceTime.IsTimeSliceCurve, MorganTianLib.GeneralizedSpaceTime.HorizontalMetric.horizontalCurveSpeed, MorganTianLib.GeneralizedSpaceTime.HorizontalMetric.horizontalCurveLength, MorganTianLib.GeneralizedSpaceTime.HorizontalMetric.horizontalCurveSpeed_constScale, MorganTianLib.GeneralizedSpaceTime.HorizontalMetric.horizontalCurveLength_constScale, MorganTianLib.GeneralizedSpaceTime.HorizontalMetric.timeSliceEDist, MorganTianLib.GeneralizedSpaceTime.HorizontalMetric.timeSliceEDist_self, MorganTianLib.GeneralizedRicciFlow.timeSliceBall, MorganTianLib.GeneralizedRicciFlow.timeSliceBall_subset_timeSlice, MorganTianLib.GeneralizedRicciFlow.mem_timeSliceBall_self, MorganTianLib.GeneralizedRicciFlow.ForwardParabolicNeighborhood, MorganTianLib.GeneralizedRicciFlow.ForwardParabolicNeighborhood.image, MorganTianLib.GeneralizedRicciFlow.ForwardParabolicNeighborhood.center_mem, MorganTianLib.GeneralizedRicciFlow.ForwardParabolicNeighborhood.image_inter_centerSlice, MorganTianLib.GeneralizedRicciFlow.ForwardParabolicNeighborhood.time_image, MorganTianLib.GeneralizedRicciFlow.BackwardParabolicNeighborhood, MorganTianLib.GeneralizedRicciFlow.BackwardParabolicNeighborhood.image, MorganTianLib.GeneralizedRicciFlow.BackwardParabolicNeighborhood.center_mem, MorganTianLib.GeneralizedRicciFlow.BackwardParabolicNeighborhood.image_inter_centerSlice, MorganTianLib.GeneralizedRicciFlow.BackwardParabolicNeighborhood.time_image, MorganTianLib.GeneralizedSpaceTime.CompatibleTimeSliceEmbedding.image_eq_of_eqOn_Icc, MorganTianLib.GeneralizedSpaceTime.CompatibleTimeSliceEmbedding.image_eq_of_closure, MorganTianLib.GeneralizedSpaceTime.CompatibleTimeSliceEmbedding.image_eq_of_isClosed, MorganTianLib.GeneralizedSpaceTime.CompatibleTimeSliceEmbedding.image_eq_of_embedding, MorganTianLib.GeneralizedRicciFlow.ForwardParabolicNeighborhood.embedding_center_eq, MorganTianLib.GeneralizedRicciFlow.ForwardParabolicNeighborhood.image_inter_centerSlice_eq, MorganTianLib.GeneralizedRicciFlow.ForwardParabolicNeighborhood.image_eq_of_eqOn_Icc, MorganTianLib.GeneralizedRicciFlow.BackwardParabolicNeighborhood.embedding_center_eq, MorganTianLib.GeneralizedRicciFlow.BackwardParabolicNeighborhood.image_inter_centerSlice_eq, MorganTianLib.GeneralizedRicciFlow.BackwardParabolicNeighborhood.image_eq_of_eqOn_Icc}

Let $({\mathcal M},G)$ be a generalized Ricci flow and let $x$ be a
point of space-time. Set $t={\bf t}(x)$. For any  $r>0$ we define
$B(x,t,r)\subset M_t$ to be the metric ball of radius $r$ centered
at $x$ in the Riemannian manifold $(M_t,g(t))$. For any $\Delta t>0$
we say that $P(x,t,r,\Delta t)$,
respectively, $P(x,r,t,-\Delta
t)$, {\em exists} in ${\mathcal M}$ if there is an embedding
$B(x,t,r)\times [t,t+\Delta t]$, respectively, $B(x,t,r)\times
[t-\Delta t,t]$, into ${\mathcal M}$ compatible with time and the
vector field. When this embedding exists, its image is defined to be the {\em
forward parabolic neighborhood}
$P(x,t,r,\Delta t)$, respectively the {\em backward parabolic
neighborhood} $P(x,t,r,-\Delta t)$.
See the figure in [Morgan--Tian].
\end{definition}
